\subsection*{19AZ$^{(6g)}$. Post-closure normalization of the live OP~1 / OP~5 Step~(B) burden labels}
\label{sec:post-closure-normalization-op1-op5b-burden-labels}
After 19AZ$^{(5d)}$ and 19AZ$^{(6f)}$, the two burdens that had governed the immediate post-release continuation line are no longer honestly read as open live kernels on the retained corridors. The right normalization step is therefore not to keep repeating the old parked labels, but to distinguish sharply between \emph{retained local kernel closure} and any broader global or machine-theoretic reading that still remains conditional.

\begin{quote}\small
\noindent\textbf{Reader cue.} Read the next block as a post-closure normalization note rather than as a new theorematic layer. Its purpose is to update the active burden language after the two local kernel closures, so that later sections do not keep speaking as if OP~1 and OP~5 Step~(B) were still parked in exactly the same way as before r01--r10.
\end{quote}

\begin{definition}[Post-closure burden labels for the retained OP kernels]
\label{def:post-closure-burden-labels-retained-op-kernels}
On the present working line we use the following normalized burden labels.
\begin{enumerate}[leftmargin=2em]
  \item \textbf{OP~1 retained local kernel.} This label denotes the OP~1 corridor problem on the intended admissible windows together with the standing retained no-second-carrier and strict-upper-gap exclusions. After Theorem~\ref{thm:closure-live-op1-kernel-standing-retained-exclusion-disciplines}, this label is read as \emph{locally discharged}.
  \item \textbf{OP~5 Step~(B) retained local kernel.} This label denotes the routed branch-shell corridor problem together with the standing retained branch-shell and shell-transfer disciplines. After Corollary~\ref{cor:live-op5b-kernel-reduces-to-one-common-shell-stable-carrier-package-deficit-problem}, this label is read as \emph{locally discharged}.
  \item \textbf{Broader conditional machine-theoretic reading.} This label denotes stronger global closure or interpretation language beyond the retained corridors and certified routes. It remains \emph{conditional} unless reopened and proved on its own terms.
\end{enumerate}
\end{definition}

\begin{proposition}[The active post-release burden language must distinguish local kernel closure from broader conditional readings]
\label{prop:active-post-release-burden-language-distinguishes-local-kernel-closure-from-broader-conditional-readings}
Under the standing hypotheses of Theorem~\ref{thm:closure-live-op1-kernel-standing-retained-exclusion-disciplines} and Corollary~\ref{cor:live-op5b-kernel-reduces-to-one-common-shell-stable-carrier-package-deficit-problem}, the following is the correct active reading of the present working line.
\begin{enumerate}[leftmargin=2em,label=(\roman*)]
  \item The OP~1 retained local kernel is no longer a parked live burden on the intended admissible windows; it is locally discharged.
  \item The OP~5 Step~(B) retained local kernel is no longer a parked live burden on the retained branch-shell corridor; it is locally discharged.
  \item Any stronger language extending these local kernel closures to a broader machine-theoretic or global reading remains conditional unless separately reopened along certified routes.
\end{enumerate}
Consequently, later front/status language in the active line should no longer speak as if OP~1 and OP~5 Step~(B) were still simply the two parked live kernels of the continuation line.
\end{proposition}

\begin{proof}
Clause (i) is exactly the content of Theorem~\ref{thm:closure-live-op1-kernel-standing-retained-exclusion-disciplines}, which discharges \(\Delta_{\mathrm{CL}}\) on the intended admissible windows under the standing retained exclusions. Clause (ii) is exactly the content of Corollary~\ref{cor:live-op5b-kernel-reduces-to-one-common-shell-stable-carrier-package-deficit-problem}, which discharges \(\Delta_{\Theta R}\) on the retained branch-shell corridor under the standing shell-transfer package. Clause (iii) is only the governing caution used throughout the present line: route-certified local kernel closure is not automatically the same as an unrestricted new global machine-theoretic closure slogan. The final sentence is therefore simply the correct normalization of the active burden language after r01--r10.
\end{proof}

\begin{remark}[What this post-closure normalization buys]
\label{rem:what-post-closure-normalization-buys}
This block does not add a third proof. It prevents a quieter but real source of drift: carrying forward the old parked-burden language after the local kernel closures have already been extracted. The active line can now speak more honestly. The tool/QAM/decoder layer is fixed infrastructure, the retained OP~1 and OP~5 Step~(B) kernels are locally discharged, and only broader conditional machine-theoretic extensions remain open for later controlled reopening.
\end{remark}


\subsection*{19AZ$^{(6h)}$. Joint post-closure synthesis and certified reopening map after the retained OP kernels}
\label{sec:joint-post-closure-synthesis-certified-reopening-map-retained-op-kernels}
After 19AZ$^{(6g)}$, the honest continuation question is no longer whether the retained OP~1 and OP~5 Step~(B) kernels should simply be reopened in their old parked form. Those kernels are already locally discharged on their certified corridors. What remains live is only the controlled extension problem: which broader tasks still survive once the retained admissible-window and retained branch-shell closures have been extracted? The next normalized step is therefore a joint synthesis block that separates \emph{closed retained kernels}, \emph{corridor-extension tasks}, and \emph{broader machine-theoretic promotion tasks}.

\begin{quote}\small
\noindent\textbf{Reader cue.} Read the next block as a post-closure continuation map rather than as a new local closure proof. Its purpose is to say what kinds of work are still honestly open after the retained OP~1 and OP~5 Step~(B) kernels have been discharged, so that later continuation notes do not quietly reopen the wrong burden.
\end{quote}

\begin{definition}[Certified post-closure task classes after the retained OP kernel closures]
\label{def:certified-post-closure-task-classes-after-retained-op-kernel-closures}
Under the standing hypotheses of Theorem~\ref{thm:closure-live-op1-kernel-standing-retained-exclusion-disciplines} and Corollary~\ref{cor:closure-live-op5b-kernel-under-standing-retained-branch-shell-and-shell-transfer-disciplines}, we distinguish the following active task classes.
\begin{enumerate}[leftmargin=2em]
  \item \textbf{Closed retained kernels.} These are the retained OP~1 admissible-window kernel and the retained OP~5 Step~(B) branch-shell kernel. They are not honest reopening targets unless one explicitly drops one of the standing retained disciplines.
  \item \textbf{Corridor-extension tasks.} These ask whether the local closure statements persist beyond the presently certified retained corridors --- for OP~1 beyond the intended admissible-window package, and for OP~5 Step~(B) beyond the retained branch-shell/shell-transfer corridor.
  \item \textbf{Broader machine-theoretic promotion tasks.} These ask whether the present tool/QAM/decoder line, together with the observer-side machine-shadow package, can be promoted from route-certified local closure language to a stronger global machine-theoretic reading.
  \item \textbf{Diagnostic task selection.} This is the use of the Monte-Carlo/fusion probe layer, residual-signature words, and decoder outputs to decide \emph{which} certified reopening route should be tested first. It is diagnostic infrastructure, not a replacement for the structural theorematic work in items~(ii)--(iii).
\end{enumerate}
\end{definition}

\begin{proposition}[After the retained kernel closures, the active continuation line consists only of extension or promotion tasks]
\label{prop:after-retained-kernel-closures-active-continuation-line-consists-only-of-extension-or-promotion-tasks}
Assume Theorem~\ref{thm:closure-live-op1-kernel-standing-retained-exclusion-disciplines}, Corollary~\ref{cor:closure-live-op5b-kernel-under-standing-retained-branch-shell-and-shell-transfer-disciplines}, and Definition~\ref{def:certified-post-closure-task-classes-after-retained-op-kernel-closures}. Then the present working line has the following normalized continuation status.
\begin{enumerate}[leftmargin=2em,label=(\roman*)]
  \item The retained OP~1 kernel is no longer an honest live burden on the intended admissible windows.
  \item The retained OP~5 Step~(B) kernel is no longer an honest live burden on the retained branch-shell corridor.
  \item Any genuine new burden must therefore be formulated either as a corridor-extension task or as a broader machine-theoretic promotion task.
  \item The Monte-Carlo/fusion probe layer serves only to classify and prioritize such tasks through the certified decoder routes; it does not by itself reopen a discharged retained kernel.
\end{enumerate}
Consequently, the active post-r11 continuation line is normalized as
\[
\text{closed retained kernels}
\ \to\ 
\text{extension or promotion tasks only},
\]
rather than as a return to the pre-closure parked-burden language.
\end{proposition}

\begin{proof}
Clauses~(i)--(ii) are exactly the retained local closure statements extracted in Theorem~\ref{thm:closure-live-op1-kernel-standing-retained-exclusion-disciplines} and Corollary~\ref{cor:closure-live-op5b-kernel-under-standing-retained-branch-shell-and-shell-transfer-disciplines}. Therefore the old reading --- ``two parked live kernels waiting to be reopened'' --- is no longer the correct active description of the present line. Once those two kernels are discharged on their certified corridors, any honest new burden must either enlarge one of the certified corridors or seek a stronger global/machine-theoretic promotion beyond them; this is exactly clauses~(iii) and the first three task classes of Definition~\ref{def:certified-post-closure-task-classes-after-retained-op-kernel-closures}. Clause~(iv) is merely the disciplined reading already built into the QAM/decoder architecture: residual-signature routing chooses which certified reopening route is tested first, but it does not substitute for the structural arguments that would be needed to prove an extension or promotion task. The displayed normalized continuation status is therefore the correct post-closure synthesis.
\end{proof}

\begin{corollary}[First normalized continuation priorities after the retained OP kernel closures]
\label{cor:first-normalized-continuation-priorities-after-retained-op-kernel-closures}
Under Proposition~\ref{prop:after-retained-kernel-closures-active-continuation-line-consists-only-of-extension-or-promotion-tasks}, the first honest continuation priorities are:
\begin{enumerate}[leftmargin=2em,label=(\alph*)]
  \item test corridor-extension statements before stronger global promotion claims,
  \item use the Monte-Carlo/fusion probe layer only as a task selector among certified routes, and
  \item keep any new reopening note explicit about which retained discipline or corridor boundary is being enlarged.
\end{enumerate}
\end{corollary}

\begin{remark}[What this joint post-closure synthesis buys]
\label{rem:what-joint-post-closure-synthesis-buys}
This block does not prove a third closure theorem. It prevents the next continuation stage from drifting back into the less honest pre-closure language. The document can now move forward with a clean separation: the retained local OP kernels are closed, the tool/QAM/decoder layer is stable infrastructure, the probe layer is diagnostic guidance, and only extension/promotion problems remain as genuinely open live tasks.
\end{remark}


\subsection*{19AZ$^{(6i)}$. First certified corridor-extension task after the retained OP kernel closures}
\label{sec:first-certified-corridor-extension-task-after-retained-op-kernel-closures}
After 19AZ$^{(6h)}$, the next honest step is not a vague reopening of the old parked burdens, and it is not yet a broad machine-theoretic promotion claim. The smallest certified continuation task is to test whether one of the already closed retained kernels survives a \emph{minimal corridor enlargement} without changing the governing local route discipline. On the present line, the cleanest first pilot is the OP~1 admissible-window side: it stays entirely inside the already extracted IP-chain, asks for no new branch-shell promotion, and is therefore the smallest honest extension task after the retained local kernel closures.

\begin{quote}\small
\noindent\textbf{Reader cue.} Read the next block as the first post-closure extension pilot, not as a third closure theorem. Its purpose is to say which corridor-extension question is small enough to be the next honest task after the retained OP~1 and OP~5 Step~(B) kernels have been discharged.
\end{quote}

\begin{definition}[One-step admissible-window extension pilot for the post-closure OP~1 line]
\label{def:one-step-admissible-window-extension-pilot-post-closure-op1-line}
Under Definition~\ref{def:certified-post-closure-task-classes-after-retained-op-kernel-closures}, the \emph{one-step admissible-window extension pilot} is the task of enlarging the intended admissible-window package of Theorem~\ref{thm:closure-live-op1-kernel-standing-retained-exclusion-disciplines} by one adjacent admissible slot while keeping fixed:
\begin{enumerate}[leftmargin=2em,label=(\roman*)]
  \item the retained total-channel profile on the already certified corridor,
  \item the retained no-second-carrier and strict-upper-gap disciplines used in 19AZ$^{(5c)}$--19AZ$^{(5d)}$,
  \item the normalized reopening route
  \[
  \mathrm{TR}_4 \Longrightarrow \mathrm{IP}_{10} \Longrightarrow \mathrm{IP}_9 \Longrightarrow \mathrm{IP}_8,
  \]
  restricted now to the first not-yet-certified adjacent admissible slot.
\end{enumerate}
A failure of this pilot is called an \emph{extension-boundary witness}. It is the first licensed witness showing that the retained OP~1 closure does not transfer unchanged to the enlarged admissible-window corridor.
\end{definition}

\begin{proposition}[The first certified extension task is the one-step OP~1 admissible-window pilot]
\label{prop:first-certified-extension-task-is-one-step-op1-admissible-window-pilot}
Assume Proposition~\ref{prop:after-retained-kernel-closures-active-continuation-line-consists-only-of-extension-or-promotion-tasks} and Definition~\ref{def:one-step-admissible-window-extension-pilot-post-closure-op1-line}. Then the first certified continuation task after the retained OP~1 / OP~5 Step~(B) kernel closures is the one-step admissible-window extension pilot, in the following precise sense.
\begin{enumerate}[leftmargin=2em,label=(\roman*)]
  \item It is a genuine corridor-extension task rather than a new local closure burden, because the retained OP~1 kernel itself has already been discharged on the intended admissible windows.
  \item It is smaller than a branch-shell extension of OP~5 Step~(B), because it stays inside the already certified IP-route and does not yet ask for a new routed branch-shell transfer argument.
  \item It is smaller than a broader machine-theoretic promotion task, because it asks only whether the already certified retained OP~1 closure transfers one slot farther along the admissible-window corridor.
  \item Therefore any honest next continuation note may either test this pilot directly or explain explicitly why a different extension/promotion task has been chosen in its place.
\end{enumerate}
\end{proposition}

\begin{proof}
Clause~(i) is immediate from Proposition~\ref{prop:after-retained-kernel-closures-active-continuation-line-consists-only-of-extension-or-promotion-tasks}: after the retained kernel closures, the old OP~1 burden is no longer live on the already certified corridor, so the only honest next OP~1 question is whether the closure transfers beyond that corridor. Clause~(ii) is purely a size comparison inside the already normalized route language. The OP~1 pilot stays on the extracted IP side and asks for no new branch/shell screening beyond what has already been certified; by contrast, an OP~5 Step~(B) corridor enlargement would immediately reopen the routed branch-shell transfer package. Clause~(iii) is again only the controlled post-closure reading already fixed in 19AZ$^{(6h)}$: a one-step corridor enlargement is strictly smaller than a global machine-theoretic promotion claim. Clause~(iv) simply records the resulting working priority. Hence the one-step admissible-window extension pilot is the first certified extension task on the present post-closure line.
\end{proof}

\begin{corollary}[Normalized immediate working priority after the retained kernel closures]
\label{cor:normalized-immediate-working-priority-after-retained-kernel-closures}
Under Proposition~\ref{prop:first-certified-extension-task-is-one-step-op1-admissible-window-pilot}, the immediate post-r12 working priority is:
\[
\text{test one-step OP~1 admissible-window transfer first}
\quad\text{before}\quad
\text{broader branch-shell extension or machine-theoretic promotion.}
\]
Any Monte-Carlo/fusion-probe use on this line is therefore read only as a selector for the first adjacent admissible slot to test; it does not itself prove the corridor transfer.
\end{corollary}

\begin{remark}[What this first certified extension pilot buys]
\label{rem:what-first-certified-extension-pilot-buys}
This block does not enlarge the certified corridor yet. It prevents a more common post-closure drift: jumping immediately from local kernel closure to a much larger machine-theoretic narrative because nothing smaller has been named. The present line now has a disciplined next target. If the OP~1 closure transfers one admissible slot farther, the corridor genuinely grows. If it fails, the failure appears as an explicit extension-boundary witness rather than as a vague reopening of the old burden language.
\end{remark}


\subsection*{19AZ$^{(6j)}$. First boundary-witness localization for the one-step OP~1 admissible-window extension pilot}
\label{sec:first-boundary-witness-localization-one-step-op1-admissible-window-extension-pilot}
After 19AZ$^{(6i)}$, the next honest step is to say where a failure of the one-step admissible-window extension pilot could still live. Because the retained OP~1 kernel has already been discharged on the certified corridor, an extension failure is not allowed to reappear as a vague shadow on the old admissible windows. If the pilot fails, the failure must localize at the first newly adjoined admissible slot and must do so through the same retained IP-route that already governs the closed corridor.

\begin{quote}\small
\noindent\textbf{Reader cue.} Read the next block as the first localization of the post-closure extension pilot, not as a second closure theorem. Its purpose is to show that any honest failure of the one-step OP~1 corridor extension must already appear at the first newly adjoined admissible slot, rather than as a hidden defect on the already certified retained corridor.
\end{quote}

\begin{definition}[First-new-slot boundary witness for the one-step admissible-window extension pilot]
\label{def:first-new-slot-boundary-witness-one-step-admissible-window-extension-pilot}
Under Definition~\ref{def:one-step-admissible-window-extension-pilot-post-closure-op1-line}, let $\sigma_{+}$ denote the first adjacent admissible slot not yet contained in the already certified retained OP~1 corridor. A \emph{first-new-slot boundary witness} is a witness obstructing the one-step admissible-window extension pilot such that:
\begin{enumerate}[leftmargin=2em,label=(\roman*)]
  \item its restriction to the already certified retained admissible-window corridor is silent, in the sense that it does not create a new live witness on the corridor already discharged by Theorem~\ref{thm:closure-live-op1-kernel-standing-retained-exclusion-disciplines},
  \item its first nontrivial support appears at $\sigma_{+}$, and
  \item it is read through the same normalized route
  \[
  \mathrm{TR}_4 \Longrightarrow \mathrm{IP}_{10} \Longrightarrow \mathrm{IP}_9 \Longrightarrow \mathrm{IP}_8,
  \]
  now applied to the enlarged corridor whose only new admissible data sit at $\sigma_{+}$.
\end{enumerate}
Such a witness may be read either as a duplicate-signature type failure or as a retained-discipline break, but in both cases its first honest support lies at the newly adjoined admissible slot.
\end{definition}

\begin{proposition}[Boundary witnesses for the one-step admissible-window pilot localize at the first new slot]
\label{prop:boundary-witnesses-one-step-admissible-window-pilot-localize-first-new-slot}
Assume Definition~\ref{def:one-step-admissible-window-extension-pilot-post-closure-op1-line} and Proposition~\ref{prop:first-certified-extension-task-is-one-step-op1-admissible-window-pilot}. Then:
\begin{enumerate}[leftmargin=2em,label=(\roman*)]
  \item if the one-step admissible-window extension pilot fails, there exists a first-new-slot boundary witness in the sense of Definition~\ref{def:first-new-slot-boundary-witness-one-step-admissible-window-extension-pilot};
  \item conversely, any first-new-slot boundary witness blocks the one-step admissible-window extension pilot;
  \item therefore the one-step admissible-window extension pilot is equivalent to excluding first-new-slot boundary witnesses on the first newly adjoined admissible slot.
\end{enumerate}
\end{proposition}

\begin{proof}
Because Theorem~\ref{thm:closure-live-op1-kernel-standing-retained-exclusion-disciplines} already discharges the retained OP~1 live kernel on the certified admissible-window corridor, a failure of the one-step extension pilot cannot honestly be supported only on that old corridor: such a witness would merely be a re-labeling of a burden already closed. Hence any genuine failure of the enlarged-corridor pilot must first become nontrivial at the first admissible slot that was not yet part of the certified corridor, namely $\sigma_{+}$. This gives clause~(i). Clause~(ii) is immediate from the definition: a witness whose first nontrivial support appears at the newly adjoined admissible slot is exactly a witness preventing the closure statement from transferring one slot farther. Clause~(iii) is just the normalization of (i)--(ii): after the retained kernel closure, the one-step extension task is neither a global reopening nor a diffuse corridor failure, but precisely the exclusion of first-new-slot boundary witnesses routed through the already extracted IP-chain.
\end{proof}

\begin{corollary}[Normalized immediate target for the one-step OP~1 extension pilot]
\label{cor:normalized-immediate-target-one-step-op1-extension-pilot}
Under Proposition~\ref{prop:boundary-witnesses-one-step-admissible-window-pilot-localize-first-new-slot}, the immediate post-r13 OP~1 continuation target is:
\[
\text{exclude first-new-slot boundary witnesses at } \sigma_{+}
\]
for the first newly adjoined admissible slot, rather than reopening the already certified retained OP~1 corridor. Any Monte-Carlo/fusion-probe use on this line is therefore read only as a selector for the first adjacent admissible slot and for the first candidate witness channel on that slot.
\end{corollary}

\begin{remark}[What this first extension-boundary localization buys]
\label{rem:what-first-extension-boundary-localization-buys}
This block does not yet enlarge the OP~1 corridor. It makes the next search honest and much smaller. If the one-step admissible-window transfer succeeds, the corridor grows by one slot. If it fails, the failure no longer appears as a mysterious return of the old OP~1 burden, but as a first-new-slot boundary witness supported at the newly adjoined admissible slot and routed through the already certified retained IP-chain.
\end{remark}



\subsection*{19AZ$^{(6k)}$. First channel split for first-new-slot boundary witnesses on the one-step OP~1 admissible-window extension pilot}
\label{sec:first-channel-split-first-new-slot-boundary-witnesses-one-step-op1-pilot}
After 19AZ$^{(6j)}$, the next honest question is not whether the pilot fails somewhere on the old corridor, but what kind of failure could still survive at the first newly adjoined admissible slot. Once the failure is localized at $\sigma_{+}$, there are only two structurally honest readings left: either the new slot creates a duplicate-signature type obstruction, or one of the standing retained exclusion disciplines themselves fails to transfer to the first new slot.

\begin{quote}\small
\noindent\textbf{Reader cue.} Read the next block as a channel split for the first-new-slot boundary witness, not as a new closure claim. Its purpose is to separate the extension pilot into the two only honest failure modes at the first newly adjoined admissible slot: duplicate-signature type versus retained-discipline break.
\end{quote}

\begin{definition}[Duplicate-signature type versus retained-discipline-break type at the first new admissible slot]
\label{def:duplicate-signature-vs-retained-discipline-break-first-new-slot-op1-pilot}
Assume Definition~\ref{def:first-new-slot-boundary-witness-one-step-admissible-window-extension-pilot}. Let $W_{+}$ be a first-new-slot boundary witness supported at $\sigma_{+}$.
\begin{enumerate}[leftmargin=2em,label=(\roman*)]
  \item We say that $W_{+}$ is of \emph{duplicate-signature type} if, on the enlarged corridor obtained by adjoining $\sigma_{+}$, the obstruction is still read as the appearance of two distinct admissible primitive completions with the same retained pinning signature, with the first honest nontrivial support of the duplicating configuration appearing at $\sigma_{+}$.
  \item We say that $W_{+}$ is of \emph{retained-discipline-break type} if the obstruction is already detected as a failure of one of the standing retained exclusion disciplines to transfer from the certified corridor to the first newly adjoined admissible slot $\sigma_{+}$; concretely, the witness is read as a first-slot failure of the retained no-second-carrier discipline or of the retained strict-upper-gap discipline.
\end{enumerate}
A first-new-slot boundary witness may in principle present both readings simultaneously, but no third autonomous channel is licensed on the current post-closure OP~1 line.
\end{definition}

\begin{proposition}[Every first-new-slot boundary witness falls into the duplicate-signature or retained-discipline-break channel]
\label{prop:first-new-slot-boundary-witnesses-split-duplicate-signature-vs-retained-discipline-break}
Assume Proposition~\ref{prop:boundary-witnesses-one-step-admissible-window-pilot-localize-first-new-slot}. Then every first-new-slot boundary witness $W_{+}$ belongs to at least one of the two classes in Definition~\ref{def:duplicate-signature-vs-retained-discipline-break-first-new-slot-op1-pilot}. More precisely:
\begin{enumerate}[leftmargin=2em,label=(\roman*)]
  \item if the standing retained exclusion disciplines transfer from the certified corridor to $\sigma_{+}$, then any first-new-slot boundary witness is necessarily of duplicate-signature type;
  \item if a first-new-slot boundary witness is not of duplicate-signature type, then it is necessarily of retained-discipline-break type;
  \item hence the one-step admissible-window extension pilot can fail only through a duplicate-signature witness at $\sigma_{+}$ or through a first-slot break of the retained exclusion disciplines.
\end{enumerate}
\end{proposition}

\begin{proof}
By Proposition~\ref{prop:boundary-witnesses-one-step-admissible-window-pilot-localize-first-new-slot}, any honest failure of the pilot is already localized at the first newly adjoined admissible slot $\sigma_{+}$. Once we are there, there are only two ways for the extension attempt to remain obstructed. Either the standing retained exclusion disciplines still govern the enlarged corridor up to the new slot; then the only remaining honest obstruction is that the new slot supports a duplicate-signature configuration of the same general kind already normalized in the retained kernel analysis, now first becoming nontrivial at $\sigma_{+}$. This is clause~(i). Or else the obstruction appears before any such duplicate-signature comparison can even be completed, namely because one of the standing retained exclusion disciplines itself no longer transfers to the first new slot. This is clause~(ii). Clause~(iii) is then only the normalized summary: after the retained kernel closure and after the first-new-slot localization, the extension pilot has exactly two licensed failure channels and no third autonomous residue.
\end{proof}

\begin{corollary}[Immediate post-r14 continuation target for the one-step OP~1 extension pilot]
\label{cor:immediate-post-r14-continuation-target-one-step-op1-extension-pilot}
Under Proposition~\ref{prop:first-new-slot-boundary-witnesses-split-duplicate-signature-vs-retained-discipline-break}, the immediate post-r14 continuation target splits into the following two certified checks at $\sigma_{+}$:
\[
\text{(a) test transfer of the standing retained exclusion disciplines to } \sigma_{+};
\]
\[
\text{(b) if transfer holds, exclude duplicate-signature witnesses at } \sigma_{+}.
\]
Any Monte-Carlo/fusion-probe use on this line is therefore read only as a selector between these two certified checks and not as an independent proof of either one.
\end{corollary}

\begin{remark}[What this first channel split at the new slot buys]
\label{rem:what-first-channel-split-at-new-slot-buys}
This block does not yet extend the corridor by one slot. It tells us exactly what must be checked next. The search is no longer allowed to wander over the old retained corridor or over vague extension rhetoric. After r14, a one-step OP~1 extension failure can only mean one of two things at the first newly adjoined admissible slot: either the retained exclusion disciplines themselves do not transfer there, or else they do transfer and one must exclude a new duplicate-signature witness supported at $\sigma_{+}$.
\end{remark}





\subsection*{19AZ$^{(6l)}$. Transfer of the standing retained exclusion disciplines to the first new admissible slot on the one-step OP~1 extension pilot}
\label{sec:transfer-standing-retained-exclusion-disciplines-first-new-slot-one-step-op1-pilot}
After 19AZ$^{(6k)}$, the next honest question is whether the standing retained exclusion disciplines themselves actually survive the one-step enlargement to the first newly adjoined admissible slot $\sigma_{+}$. On the present line this is the smallest structural transfer test: if the one-step pilot really stays on the same retained IP-route, keeps the same retained total-channel profile and optimum packet, and asks for no new branch/shell promotion before the new slot is read, then the no-second-carrier and strict-upper-gap exclusions should not suddenly fail at $\sigma_{+}$. The next step is therefore to state this transfer in the smallest certified form and thereby remove the retained-discipline-break channel whenever the pilot remains a genuine one-step IP-side enlargement.

\begin{quote}\small
\noindent\textbf{Reader cue.} Read the next block as a first-slot transfer theorem for the standing retained exclusion disciplines, not as the corridor extension itself. Its purpose is to remove the retained-discipline-break channel at $\sigma_{+}$ whenever the one-step pilot truly remains inside the same retained IP-side comparison packet.
\end{quote}

\begin{definition}[First-slot retained-discipline transfer certificate for the one-step OP~1 pilot]
\label{def:first-slot-retained-discipline-transfer-certificate-one-step-op1-pilot}
Under Definition~\ref{def:one-step-admissible-window-extension-pilot-post-closure-op1-line}, we say that the one-step OP~1 admissible-window extension pilot carries the \emph{first-slot retained-discipline transfer certificate} if all of the following hold on the enlarged corridor through the first newly adjoined admissible slot $\sigma_{+}$:
\begin{enumerate}[leftmargin=2em,label=(\roman*)]
  \item the extension is still read through the same retained IP-route \(\mathrm{TR}_4 \Rightarrow \mathrm{IP}_{10} \Rightarrow \mathrm{IP}_9 \Rightarrow \mathrm{IP}_8\) used on the certified corridor;
  \item the retained total-channel profile and retained optimum packet used in the closed kernel analysis remain unchanged up to and including $\sigma_{+}$;
  \item no new branch/shell promotion or BS/BL-side reopening is licensed before the first new slot is read;
  \item every candidate primitive occupancy or admissible completion involving $\sigma_{+}$ is still compared against the same retained comparison packet that governed the certified corridor.
\end{enumerate}
In this situation we say that the standing retained no-second-carrier and strict-upper-gap disciplines \emph{transfer to the first new admissible slot}.
\end{definition}

\begin{proposition}[Under first-slot transfer, the retained-discipline-break channel at $\sigma_{+}$ is absent]
\label{prop:first-slot-transfer-removes-retained-discipline-break-channel-sigma-plus}
Assume Proposition~\ref{prop:first-new-slot-boundary-witnesses-split-duplicate-signature-vs-retained-discipline-break} and Definition~\ref{def:first-slot-retained-discipline-transfer-certificate-one-step-op1-pilot}. Then no first-new-slot boundary witness supported at $\sigma_{+}$ can be of retained-discipline-break type. Equivalently, whenever the one-step extension pilot carries the first-slot retained-discipline transfer certificate, the standing retained no-second-carrier and strict-upper-gap exclusions extend from the certified corridor to $\sigma_{+}$.
\end{proposition}

\begin{proof}
A retained-discipline-break witness at $\sigma_{+}$ would mean that, before any genuine duplicate-signature comparison has become the first honest obstruction, one of the two standing retained exclusion disciplines already fails to survive the one-step enlargement. But the transfer certificate says exactly that the pilot has not changed the governing route, has not changed the retained total-channel profile or optimum packet, has not introduced a new branch/shell promotion before reading the new slot, and still compares every candidate occupancy/completion involving $\sigma_{+}$ against the same retained comparison packet used on the closed corridor. Under those hypotheses there is no new structural mechanism available that could create a first-slot failure of the retained no-second-carrier discipline or of the retained strict-upper-gap discipline without first leaving the certified one-step IP-side enlargement setting itself. Therefore the retained-discipline-break channel is absent at $\sigma_{+}$.
\end{proof}

\begin{corollary}[Under first-slot transfer, any surviving boundary witness at $\sigma_{+}$ is of duplicate-signature type]
\label{cor:under-first-slot-transfer-boundary-witness-at-sigma-plus-is-duplicate-signature-type}
Assume Definition~\ref{def:first-slot-retained-discipline-transfer-certificate-one-step-op1-pilot}. Then any first-new-slot boundary witness supported at $\sigma_{+}$ is necessarily of duplicate-signature type. In particular, under first-slot transfer the one-step admissible-window extension pilot can fail only through a new duplicate-signature witness at the first newly adjoined admissible slot.
\end{corollary}

\begin{remark}[What the first-slot retained-discipline transfer buys]
\label{rem:what-first-slot-retained-discipline-transfer-buys}
This does not yet prove the one-step extension pilot. It removes one entire failure channel. After the present step, the only honest remaining obstruction at $\sigma_{+}$ is the duplicate-signature channel, provided the pilot truly stays inside the same retained IP-side comparison packet. This is exactly the kind of reduction for which the Monte-Carlo/fusion-probe layer is useful: it does not prove the transfer theorem, but it can help tell us early whether a candidate run still looks like a genuine one-step IP-side enlargement or whether it is already drifting into a branch/shell promotion where the certificate would no longer apply.
\end{remark}


\noindent\textbf{Reader cue.} Read the next block as the final test on the one-step OP~1 admissible-window extension pilot. Its purpose is to decide whether the only remaining first-new-slot boundary channel --- duplicate-signature at $\sigma_{+}$ --- can still survive once the standing retained exclusion disciplines have been transferred to that new slot. If it cannot, then the one-step pilot itself is certified.

\subsection*{19AZ$^{(6m)}$. Duplicate-signature witnesses at the first new admissible slot collapse under the transferred retained exclusion disciplines}
\label{sec:duplicate-signature-witnesses-first-new-slot-collapse-under-transferred-retained-exclusion-disciplines}
After 19AZ$^{(6l)}$, the one-step OP~1 admissible-window extension pilot has only one honest remaining failure channel: a duplicate-signature witness supported at the first newly adjoined admissible slot $\sigma_{+}$. The next honest question is therefore no longer whether the retained exclusion disciplines survive there --- that was the content of the transfer certificate --- but whether the same duplicate-signature contradiction package that closed the retained OP~1 kernel also closes the first-new-slot case once those disciplines have been transferred. On the present line the answer is yes: under first-slot transfer, any surviving duplicate-signature configuration at $\sigma_{+}$ is again forced into the same occupancy-side versus completion-side split, and each side is discharged by the transferred retained no-second-carrier / strict-upper-gap exclusions.

\begin{definition}[First-new-slot duplicate-signature witness on the one-step OP~1 extension pilot]
\label{def:first-new-slot-duplicate-signature-witness-one-step-op1-extension-pilot}
Under Definition~\ref{def:one-step-admissible-window-extension-pilot-post-closure-op1-line}, a boundary witness supported at the first newly adjoined admissible slot $\sigma_{+}$ is called a \emph{first-new-slot duplicate-signature witness} if it is of duplicate-signature type in the sense of Proposition~\ref{prop:first-new-slot-boundary-witnesses-split-duplicate-signature-vs-retained-discipline-break}. Equivalently, it is a boundary witness at $\sigma_{+}$ for which the standing retained exclusion disciplines transfer to the new slot, but the enlarged corridor still carries two primitive loop/completion candidates with the same retained pinning signature when the first new slot is read through the same retained comparison packet.
\end{definition}

\begin{proposition}[Under first-slot transfer, no duplicate-signature witness survives at $\sigma_{+}$]
\label{prop:under-first-slot-transfer-no-duplicate-signature-witness-survives-at-sigma-plus}
Assume Definition~\ref{def:first-slot-retained-discipline-transfer-certificate-one-step-op1-pilot}, Definition~\ref{def:first-new-slot-duplicate-signature-witness-one-step-op1-extension-pilot}, Proposition~\ref{prop:occupancy-side-duplicate-signature-witnesses-excluded-retained-op1-corridor}, and Proposition~\ref{prop:completion-side-duplicate-signature-witnesses-excluded-retained-op1-corridor}. Then no first-new-slot duplicate-signature witness supported at $\sigma_{+}$ exists.
\end{proposition}

\begin{proof}
Assume for contradiction that a first-new-slot duplicate-signature witness exists at $\sigma_{+}$. By definition, the one-step pilot still carries the same retained IP-route, retained total-channel profile, retained optimum packet, and retained comparison packet up to and including the first new slot. Therefore the putative duplicate-signature configuration at $\sigma_{+}$ is of the same structural species as in the retained OP~1 kernel analysis: once it is read through the common retained comparison packet, it again splits into an occupancy-side versus completion-side contradiction channel exactly as in Proposition~\ref{prop:duplicate-signature-witness-feeds-standing-op1-exclusion-gate}. Because the first-slot retained-discipline transfer certificate is in force, the retained no-second-carrier and retained strict-upper-gap exclusions have already been transferred to the enlarged corridor through $\sigma_{+}$. Proposition~\ref{prop:occupancy-side-duplicate-signature-witnesses-excluded-retained-op1-corridor} therefore removes the occupancy-side channel, while Proposition~\ref{prop:completion-side-duplicate-signature-witnesses-excluded-retained-op1-corridor} removes the completion-side channel. Since no third duplicate-signature channel is licensed, the witness cannot exist. This contradiction proves the claim.
\end{proof}

\begin{corollary}[Closure of the one-step OP~1 admissible-window extension pilot under first-slot transfer]
\label{cor:closure-one-step-op1-admissible-window-extension-pilot-under-first-slot-transfer}
Assume Definition~\ref{def:first-slot-retained-discipline-transfer-certificate-one-step-op1-pilot}. Then the one-step OP~1 admissible-window extension pilot closes. Equivalently, under first-slot transfer there is no first-new-slot boundary witness at $\sigma_{+}$, so the retained OP~1 corridor extends by one admissible slot without opening any new live burden.
\end{corollary}

\begin{proof}
By Corollary~\ref{cor:under-first-slot-transfer-boundary-witness-at-sigma-plus-is-duplicate-signature-type}, any surviving first-new-slot boundary witness at $\sigma_{+}$ would have to be of duplicate-signature type. Proposition~\ref{prop:under-first-slot-transfer-no-duplicate-signature-witness-survives-at-sigma-plus} excludes exactly that possibility. Hence no first-new-slot boundary witness survives, and the one-step admissible-window extension pilot closes.
\end{proof}

\begin{remark}[What the first certified one-step OP~1 extension buys]
\label{rem:what-first-certified-one-step-op1-extension-buys}
This is the first actual corridor-extension result after the retained OP~1 kernel closure. It does not claim an arbitrary global OP~1 promotion. It says something much smaller and therefore much cleaner: if the first newly adjoined admissible slot really remains inside the same retained IP-side comparison packet, then the earlier duplicate-signature contradiction package propagates one step further and the pilot succeeds. In particular, any genuinely new OP~1 burden beyond the retained corridor must now appear only \emph{after} this first certified slot extension or by leaving the first-slot transfer certificate itself.
\end{remark}


\noindent\textbf{Reader cue.} Read the next block as the controlled iteration schema after the first certified one-step OP~1 admissible-window extension. Its purpose is not to announce an arbitrary global promotion, but to normalize the smallest slot-by-slot rule now available: whenever the next newly adjoined admissible slot again carries the same transferred retained comparison discipline, the corridor extends by one more slot; if not, the first honest extension-boundary witness must localize exactly at the first slot where that transfer certificate fails.

\subsection*{19AZ$^{(6n)}$. First iterable slotwise extension schema after the certified one-step OP~1 admissible-window transfer}
\label{sec:first-iterable-slotwise-extension-schema-after-certified-one-step-op1-admissible-window-transfer}
After 19AZ$^{(6m)}$, the current OP~1 line has something strictly stronger than a single isolated pilot success but strictly weaker than an unrestricted global promotion: it has one certified slot-transfer mechanism. The next honest step is therefore to package this mechanism into an iterable schema. The idea is conservative. One does not claim that arbitrarily many future slots automatically close; one says only that once a current corridor has already been certified, each further admissible slot may be adjoined by the same one-step argument provided the same slotwise transfer certificate is re-established at that new leading edge. In this way extension ceases to be an opaque reopening and becomes a controlled inductive procedure with a sharply localized first failure slot.

\begin{definition}[Certified slotwise OP~1 extension chain]
\label{def:certified-slotwise-op1-extension-chain}
Let $\mathcal C_0$ denote a currently certified OP~1 admissible corridor, beginning with the retained corridor already closed by Corollary~\ref{cor:closure-one-step-op1-admissible-window-extension-pilot-under-first-slot-transfer} after one certified enlargement. For $j\ge 1$, let $\sigma_j$ denote the next newly adjoined admissible slot beyond $\mathcal C_{j-1}$ and let $\mathcal C_j:=\mathcal C_{j-1}\cup\{\sigma_j\}$. We call
\[
  \mathcal C_0\subset \mathcal C_1\subset \cdots \subset \mathcal C_k
\]
a \emph{certified slotwise OP~1 extension chain of length $k$} if, for each $1\le j\le k$, the step from $\mathcal C_{j-1}$ to $\mathcal C_j$ admits the same kind of slotwise retained-discipline transfer certificate as in Definition~\ref{def:first-slot-retained-discipline-transfer-certificate-one-step-op1-pilot}, now read relative to the already certified corridor $\mathcal C_{j-1}$ and its current leading slot $\sigma_j$.
\end{definition}

\begin{proposition}[Slotwise transfer certificates iterate the one-step OP~1 extension mechanism]
\label{prop:slotwise-transfer-certificates-iterate-one-step-op1-extension-mechanism}
Assume Definition~\ref{def:certified-slotwise-op1-extension-chain}. Then every step $\mathcal C_{j-1}\leadsto \mathcal C_j$ in the chain closes, and hence the terminal corridor $\mathcal C_k$ is certified without opening a new live OP~1 burden.
\end{proposition}

\begin{proof}
Proceed by induction on $j$. The case $j=1$ is exactly the certified one-step extension mechanism already established in Corollary~\ref{cor:closure-one-step-op1-admissible-window-extension-pilot-under-first-slot-transfer}. Assume now that $\mathcal C_{j-1}$ is already certified. By definition, the next slot $\sigma_j$ comes equipped with the same slotwise retained-discipline transfer certificate, now read relative to the already certified corridor $\mathcal C_{j-1}$. Thus the step from $\mathcal C_{j-1}$ to $\mathcal C_j$ is formally the same one-step situation as before: the standing retained exclusions transfer to the new leading slot, the only honest first-new-slot boundary channel is duplicate-signature type, and that channel again collapses under the transferred retained no-second-carrier and strict-upper-gap package. Therefore $\mathcal C_j$ is certified. Induction yields the claim for every $j\le k$ and in particular for $\mathcal C_k$.
\end{proof}

\begin{corollary}[First failure of a slotwise extension attempt localizes at the first uncertified leading slot]
\label{cor:first-failure-slotwise-extension-attempt-localizes-first-uncertified-leading-slot}
Let a slotwise OP~1 extension attempt proceed successively through admissible slots beyond a certified corridor. If the attempt fails, then there is a smallest index $j_\ast$ such that the step to $\sigma_{j_\ast}$ does not admit a slotwise retained-discipline transfer certificate. Equivalently, as long as the next leading slot continues to satisfy the transferred retained comparison discipline, no new live OP~1 burden appears.
\end{corollary}

\begin{proof}
If every step admitted the slotwise retained-discipline transfer certificate, Proposition~\ref{prop:slotwise-transfer-certificates-iterate-one-step-op1-extension-mechanism} would certify the entire chain. Therefore any failure must occur at the first index where this certificate is no longer available. Minimality of that index is automatic.
\end{proof}

\begin{remark}[What the iterable slotwise extension schema now buys]
\label{rem:what-iterable-slotwise-extension-schema-now-buys}
This is the clean post-r17 upgrade of the first certified slot extension. The text still does not claim that arbitrary future admissible slots are automatically harmless. It claims something smaller and more useful: extension has become an explicit iterative schema. The Monte-Carlo/fusion-probe layer is therefore no longer searching blindly for ``whether OP~1 reopens somewhere''; it is searching for the first leading slot at which the slotwise retained-discipline transfer certificate fails or begins to look unstable. That is exactly the kind of sharply localized extension-boundary information that a routing/probe layer can accelerate without replacing the structural proof itself.
\end{remark}


\noindent\textbf{Reader cue.} Read the next block as the first concrete two-step realization of the slotwise OP~1 extension schema. Its purpose is not to announce an unrestricted corridor promotion, but to show that the new inductive format already yields a genuine $k=2$ certified chain once the second new leading slot carries the same transferred retained comparison discipline as the first.

\subsection*{19AZ$^{(6o)}$. First certified two-slot OP~1 admissible-window extension chain under repeated leading-slot transfer}
\label{sec:first-certified-two-slot-op1-admissible-window-extension-chain-under-repeated-leading-slot-transfer}
After 19AZ$^{(6n)}$, the honest question is whether the iterable schema does anything beyond packaging the already certified first new slot. The smallest nontrivial test is therefore the second newly adjoined admissible slot. One keeps the already enlarged corridor $\mathcal C_1$ fixed, adjoins the next admissible slot $\sigma_2$, and asks whether the same leading-slot transfer certificate can be re-established there. If it can, then the abstract schema immediately becomes a first genuine two-step certified chain rather than a merely formal induction template.

\begin{definition}[Second-slot transfer certificate for the first certified two-step OP~1 chain]
\label{def:second-slot-transfer-certificate-first-certified-two-step-op1-chain}
Let $\mathcal C_1$ be the corridor already certified by the one-step OP~1 admissible-window extension result of Corollary~\ref{cor:closure-one-step-op1-admissible-window-extension-pilot-under-first-slot-transfer}. Let $\sigma_2$ denote the next admissible slot beyond the leading edge of $\mathcal C_1$, and set $\mathcal C_2:=\mathcal C_1\cup\{\sigma_2\}$. We say that $\sigma_2$ carries the \emph{second-slot transfer certificate} if the step $\mathcal C_1\leadsto \mathcal C_2$ satisfies the same slotwise retained-discipline transfer conditions as in Definition~\ref{def:certified-slotwise-op1-extension-chain}, now read relative to the already certified corridor $\mathcal C_1$ and its new leading slot $\sigma_2$.
\end{definition}

\begin{proposition}[The second newly adjoined admissible slot closes under the repeated leading-slot transfer certificate]
\label{prop:second-newly-adjoined-admissible-slot-closes-under-repeated-leading-slot-transfer-certificate}
Assume Definition~\ref{def:second-slot-transfer-certificate-first-certified-two-step-op1-chain}. Then the extension step $\mathcal C_1\leadsto \mathcal C_2$ closes, and $\mathcal C_2$ is a certified OP~1 admissible corridor.
\end{proposition}

\begin{proof}
The corridor $\mathcal C_1$ is already certified by Corollary~\ref{cor:closure-one-step-op1-admissible-window-extension-pilot-under-first-slot-transfer}. By Definition~\ref{def:second-slot-transfer-certificate-first-certified-two-step-op1-chain}, the next leading slot $\sigma_2$ again satisfies exactly the slotwise retained-discipline transfer assumptions required in the iterable extension schema of Proposition~\ref{prop:slotwise-transfer-certificates-iterate-one-step-op1-extension-mechanism}. Applying that proposition with $k=2$ yields closure of the step $\mathcal C_1\leadsto \mathcal C_2$ and certifies the enlarged corridor $\mathcal C_2$.
\end{proof}

\begin{corollary}[First realized certified two-slot OP~1 admissible-window chain]
\label{cor:first-realized-certified-two-slot-op1-admissible-window-chain}
Under the repeated leading-slot transfer certificate, the current OP~1 line already contains a realized certified chain
\[
  \mathcal C_0\subset \mathcal C_1\subset \mathcal C_2,
\]
of length $2$. Equivalently, beyond the retained live-kernel closure, two consecutive admissible-slot enlargements are now certified without reopening a live OP~1 burden.
\end{corollary}

\begin{proof}
The first step $\mathcal C_0\leadsto \mathcal C_1$ is certified by Corollary~\ref{cor:closure-one-step-op1-admissible-window-extension-pilot-under-first-slot-transfer}. Proposition~\ref{prop:second-newly-adjoined-admissible-slot-closes-under-repeated-leading-slot-transfer-certificate} certifies the second step $\mathcal C_1\leadsto \mathcal C_2$. Hence the displayed chain of length $2$ is certified.
\end{proof}

\begin{remark}[What the first realized two-slot chain now buys]
\label{rem:what-first-realized-two-slot-chain-now-buys}
This is the first point at which the iterable extension schema becomes visibly more than a packaging convenience. The text still does not claim a fully promoted OP~1 corridor. It claims something smaller and already operational: once the transferred retained comparison discipline persists across two successive leading slots, the corridor extends twice without reopening the local OP~1 burden. In particular, any future extension failure must now occur only at the first admissible slot \emph{after} $\mathcal C_2$ where the leading-slot transfer certificate ceases to hold or begins to destabilize.
\end{remark}


\noindent\textbf{Reader cue.} Read the next block as a compact certification format for a finite OP~1 extension run. Its purpose is not to claim unrestricted corridor promotion, but to package a bounded number of already justified leading-slot transfers into one reusable certificate whose failure horizon is explicitly the first slot beyond the certified run.

\subsection*{19AZ$^{(6p)}$. First bounded $k$-step certificate format for finite OP~1 admissible-window extension runs}
\label{sec:first-bounded-k-step-certificate-format-for-finite-op1-admissible-window-extension-runs}
After 19AZ$^{(6o)}$, the present line has more than one isolated extension success but still less than an unrestricted promoted corridor. The right packaging is therefore a bounded certificate format: one records a finite certified run together with the leading admissible slots and their slotwise transfer certificates, and one treats the first uncertified slot beyond that run as the first honest location where extension may fail again.

\begin{definition}[Bounded $k$-step OP~1 extension certificate]
\label{def:bounded-k-step-op1-extension-certificate}
Let $k\ge 1$. A \emph{bounded $k$-step OP~1 extension certificate} is data
\[
  \mathbf K_k=(\mathcal C_0\subset\mathcal C_1\subset\cdots\subset\mathcal C_k;\,\sigma_1,\dots,\sigma_k),
\]
where for each $j=1,\dots,k$ the slot $\sigma_j$ is the newly adjoined leading admissible slot in the extension step $\mathcal C_{j-1}\leadsto \mathcal C_j$, and $\sigma_j$ carries the slotwise retained-discipline transfer certificate required by Definition~\ref{def:certified-slotwise-op1-extension-chain} relative to the already certified corridor $\mathcal C_{j-1}$.
\end{definition}

\begin{proposition}[Bounded $k$-step certificates certify the whole finite OP~1 run]
\label{prop:bounded-k-step-certificates-certify-whole-finite-op1-run}
Assume Definition~\ref{def:bounded-k-step-op1-extension-certificate}. Then the terminal corridor $\mathcal C_k$ is certified, and every intermediate step $\mathcal C_{j-1}\leadsto \mathcal C_j$ closes without reopening a live OP~1 burden.
\end{proposition}

\begin{proof}
This is exactly the repeated one-step mechanism of Proposition~\ref{prop:slotwise-transfer-certificates-iterate-one-step-op1-extension-mechanism}, now written in bounded packaged form. Since each newly adjoined leading slot $\sigma_j$ carries the required slotwise retained-discipline transfer certificate, every step closes in turn. Hence the entire finite chain up to $\mathcal C_k$ is certified.
\end{proof}

\begin{corollary}[First honest failure lies beyond a bounded certificate horizon]
\label{cor:first-honest-failure-lies-beyond-bounded-certificate-horizon}
Assume Definition~\ref{def:bounded-k-step-op1-extension-certificate}. Then any honest future extension failure must first appear at the earliest admissible slot strictly beyond the leading edge of $\mathcal C_k$ for which the next slotwise transfer certificate is absent or unstable.
\end{corollary}

\begin{proof}
Every step through $\mathcal C_k$ is already certified by Proposition~\ref{prop:bounded-k-step-certificates-certify-whole-finite-op1-run}. Therefore no honest failure can occur inside that bounded run. The first possible extension failure must lie at the first later leading slot for which the next certificate is not available.
\end{proof}

\begin{remark}[What the bounded certificate format now buys]
\label{rem:what-bounded-certificate-format-now-buys}
This is the finite-horizon normalization that the post-r17 extension line needed. Instead of speaking vaguely about an enlarging corridor, the text can now speak about a specific certified horizon $\mathbf K_k$. The Monte-Carlo/fusion-probe layer is thereby given a precise target: it need only search for the first later slot beyond $\mathbf K_k$ where the transfer certificate begins to fail, rather than resurveying the already certified run.
\end{remark}


\noindent\textbf{Reader cue.} Read the next block as the first explicit named bounded certificate already available on the present line. Its purpose is to turn the abstract finite-horizon format into a concrete certified object that can be cited directly in later routing and probe discussions.

\subsection*{19AZ$^{(6q)}$. First explicit bounded certificate $\mathbf K_2$ on the present OP~1 extension line}
\label{sec:first-explicit-bounded-certificate-k2-on-present-op1-extension-line}
After 19AZ$^{(6p)}$, the next honest step is to name the first concrete bounded certificate already realized by the current line. The smallest nontrivial such object is the two-step extension chain from 19AZ$^{(6o)}$, now repackaged as a single finite certificate. This gives the text an explicit certified horizon that later extension tests and Monte-Carlo/probe diagnostics can cite without reopening the step-by-step derivation each time.

\begin{definition}[The present two-step bounded certificate $\mathbf K_2^{\mathrm{cur}}$]
\label{def:present-two-step-bounded-certificate-k2cur}
Let
\[
  \mathbf K_2^{\mathrm{cur}}:=(\mathcal C_0\subset \mathcal C_1\subset \mathcal C_2;\,\sigma_1,\sigma_2),
\]
where $\mathcal C_0$ is the retained certified OP~1 corridor, $\sigma_1$ is the first newly adjoined admissible slot certified in Corollary~\ref{cor:closure-one-step-op1-admissible-window-extension-pilot-under-first-slot-transfer}, and $\sigma_2$ is the second newly adjoined admissible slot certified in Corollary~\ref{cor:first-realized-certified-two-slot-op1-admissible-window-chain}. We call $\mathbf K_2^{\mathrm{cur}}$ the \emph{present two-step bounded certificate}.
\end{definition}

\begin{proposition}[The current OP~1 line carries the explicit bounded certificate $\mathbf K_2^{\mathrm{cur}}$]
\label{prop:current-op1-line-carries-explicit-bounded-certificate-k2cur}
The data of Definition~\ref{def:present-two-step-bounded-certificate-k2cur} satisfy Definition~\ref{def:bounded-k-step-op1-extension-certificate} with $k=2$. Hence $\mathbf K_2^{\mathrm{cur}}$ is a valid bounded certificate on the present OP~1 line.
\end{proposition}

\begin{proof}
The step $\mathcal C_0\leadsto\mathcal C_1$ is certified by Corollary~\ref{cor:closure-one-step-op1-admissible-window-extension-pilot-under-first-slot-transfer}. The step $\mathcal C_1\leadsto\mathcal C_2$ is certified by Corollary~\ref{cor:first-realized-certified-two-slot-op1-admissible-window-chain}. Therefore the two leading slots $\sigma_1,\sigma_2$ each carry the required slotwise transfer certificate and the chain satisfies Definition~\ref{def:bounded-k-step-op1-extension-certificate} with $k=2$.
\end{proof}

\begin{corollary}[The first uncertified OP~1 slot lies strictly beyond $\mathbf K_2^{\mathrm{cur}}$]
\label{cor:first-uncertified-op1-slot-lies-strictly-beyond-k2cur}
On the present line, any honest future OP~1 admissible-window extension failure must first occur at a leading admissible slot strictly beyond the corridor $\mathcal C_2$ certified by $\mathbf K_2^{\mathrm{cur}}$.
\end{corollary}

\begin{proof}
This is the case $k=2$ of Corollary~\ref{cor:first-honest-failure-lies-beyond-bounded-certificate-horizon}, applied to Proposition~\ref{prop:current-op1-line-carries-explicit-bounded-certificate-k2cur}.
\end{proof}

\begin{remark}[What the explicit certificate $\mathbf K_2^{\mathrm{cur}}$ now buys]
\label{rem:what-explicit-certificate-k2cur-now-buys}
The current extension line now has a named finite horizon. This is useful both mathematically and diagnostically. Mathematically, later corridor-extension arguments can cite $\mathbf K_2^{\mathrm{cur}}$ as a certified base segment rather than replaying the first two transferred slots. Diagnostically, the Monte-Carlo/fusion-probe layer can now ask a sharply bounded question: where is the first later slot beyond $\mathbf K_2^{\mathrm{cur}}$ at which the next leading-slot transfer certificate destabilizes?
\end{remark}


\noindent\textbf{Reader cue.} Read the next block as the first controlled test beyond the current certified horizon $\mathbf K_2^{\mathrm{cur}}$. Its purpose is not to claim a third transferred slot automatically, but to localize the next honest extension question at the very first later leading slot and to state the exact alternative: either that slot carries the next transfer certificate and yields a certified $\mathbf K_3$, or it is the first genuine post-$\mathbf K_2^{\mathrm{cur}}$ boundary witness.

\subsection*{19AZ$^{(6r)}$. First post-$\mathbf K_2^{\mathrm{cur}}$ leading-slot test beyond the current certified horizon}
\label{sec:first-post-k2-leading-slot-test-beyond-current-certified-horizon}
After 19AZ$^{(6q)}$, the next honest step is to move from a named finite horizon to the first controlled test beyond that horizon. The current line therefore does not jump directly to an unbounded promotion claim. Instead it isolates the first later leading admissible slot beyond $\mathcal C_2$ and asks whether the same slotwise retained-discipline transfer certificate extends one more time. This is the smallest extension question that remains once $\mathbf K_2^{\mathrm{cur}}$ has been certified.

\begin{definition}[First post-$\mathbf K_2^{\mathrm{cur}}$ leading-slot test datum]
\label{def:first-post-k2-leading-slot-test-datum}
Let $\mathbf K_2^{\mathrm{cur}}=(\mathcal C_0\subset \mathcal C_1\subset \mathcal C_2;\,\sigma_1,\sigma_2)$ be the present two-step bounded certificate from Definition~\ref{def:present-two-step-bounded-certificate-k2cur}. Define $\sigma_3$ to be the first later leading admissible slot strictly beyond the certified corridor $\mathcal C_2$, whenever such a next slot is under consideration on the same retained OP~1 extension line. We call the pair $(\mathcal C_2,\sigma_3)$ the \emph{first post-$\mathbf K_2^{\mathrm{cur}}$ leading-slot test datum}.
\end{definition}

\begin{proposition}[The first honest future OP~1 extension question beyond $\mathbf K_2^{\mathrm{cur}}$ is localized at $\sigma_3$]
\label{prop:first-honest-future-op1-extension-question-beyond-k2cur-localized-at-sigma3}
Assume the setting of Definition~\ref{def:first-post-k2-leading-slot-test-datum}. Then every honest future OP~1 admissible-window extension question beyond the currently certified horizon is localized at the single leading slot $\sigma_3$. More precisely, exactly one of the following alternatives occurs:
\begin{enumerate}[label=(\roman*),leftmargin=2em]
  \item the slot $\sigma_3$ carries the same slotwise retained-discipline transfer certificate as in the previously certified steps, in which case the enlarged chain
  \[
    \mathcal C_0\subset \mathcal C_1\subset \mathcal C_2\subset \mathcal C_3
  \]
  is a valid bounded certificate $\mathbf K_3^{\mathrm{cand}}$ on the same OP~1 line;
  \item the slot $\sigma_3$ does not carry that certificate, in which case $\sigma_3$ itself is the first genuine post-$\mathbf K_2^{\mathrm{cur}}$ boundary-witness location.
\end{enumerate}
In particular, no honest future failure can occur strictly before $\sigma_3$.
\end{proposition}

\begin{proof}
By Corollary~\ref{cor:first-uncertified-op1-slot-lies-strictly-beyond-k2cur}, any honest future OP~1 extension failure must first occur at a leading admissible slot strictly beyond $\mathcal C_2$. By definition, $\sigma_3$ is the first such later leading slot under consideration. Therefore no honest future failure can occur before $\sigma_3$. If $\sigma_3$ carries the same slotwise retained-discipline transfer certificate as in the previous certified steps, then the same slotwise-extension mechanism used in Proposition~\ref{prop:bounded-k-step-certificates-certify-whole-finite-op1-run} extends the chain by one step and yields a valid bounded certificate with $k=3$. If not, then the failure is already localized at $\sigma_3$ itself. These are the only two possibilities.
\end{proof}

\begin{corollary}[The post-$\mathbf K_2^{\mathrm{cur}}$ horizon is controlled by a single next-slot test]
\label{cor:post-k2-horizon-controlled-by-single-next-slot-test}
On the present OP~1 line, the first uncertified future horizon beyond $\mathbf K_2^{\mathrm{cur}}$ is completely controlled by the single next-slot test at $\sigma_3$: either $\mathbf K_3^{\mathrm{cand}}$ is certified, or the first honest future boundary witness is already localized at $\sigma_3$.
\end{corollary}

\begin{proof}
Immediate from Proposition~\ref{prop:first-honest-future-op1-extension-question-beyond-k2cur-localized-at-sigma3}.
\end{proof}

\begin{remark}[What the first post-$\mathbf K_2^{\mathrm{cur}}$ test now buys]
\label{rem:what-first-post-k2-leading-slot-test-now-buys}
This is the first point at which the extension line leaves the already named certified horizon. The benefit is that the next search question is now maximally sharp. Neither the theorematic line nor the Monte-Carlo/fusion-probe layer needs to resurvey the already certified segment $\mathbf K_2^{\mathrm{cur}}$; both need only test whether the next leading slot $\sigma_3$ carries the same transfer certificate. If it does, a concrete $\mathbf K_3$ becomes available. If it does not, then the first genuine post-$\mathbf K_2^{\mathrm{cur}}$ boundary witness has already been found.
\end{remark}

\noindent\textbf{Reader cue.} Read the next block as the two-channel split of the first uncertified slot beyond $\mathbf K_2^{\mathrm{cur}}$. Its purpose is not yet to certify $\mathbf K_3$, but to normalize the next decision at $\sigma_3$ into the two only honest outcomes: either the same transfer certificate still holds, or the first genuine post-$\mathbf K_2^{\mathrm{cur}}$ boundary witness is already present there.

\subsection*{19AZ$^{(6s)}$. First transfer-positive versus boundary-active split at the post-$\mathbf K_2^{\mathrm{cur}}$ slot $\sigma_3$}
\label{sec:first-transfer-positive-versus-boundary-active-split-at-post-k2-slot-sigma3}
After 19AZ$^{(6r)}$, the first uncertified slot beyond the current certified horizon is no longer diffuse. It is the single leading admissible slot $\sigma_3$ beyond $\mathcal C_2$. The next honest normalization is therefore to split the status of $\sigma_3$ itself. Either $\sigma_3$ still carries the same slotwise retained-discipline transfer certificate that already certified $\sigma_1$ and $\sigma_2$, or else $\sigma_3$ is already the first active boundary location beyond $\mathbf K_2^{\mathrm{cur}}$.

\begin{definition}[Transfer-positive versus boundary-active status at $\sigma_3$]
\label{def:transfer-positive-versus-boundary-active-status-at-sigma3}
Assume the setting of Definition~\ref{def:first-post-k2-leading-slot-test-datum}. We say that the first later leading admissible slot $\sigma_3$ is
\begin{enumerate}[label=(\roman*)]
  \item \emph{transfer-positive} if it carries the same slotwise retained-discipline transfer certificate as in the certified steps leading to $\mathbf K_2^{\mathrm{cur}}$, and
  \item \emph{boundary-active} if it does not carry that certificate and therefore already hosts the first genuine post-$\mathbf K_2^{\mathrm{cur}}$ boundary witness.
\end{enumerate}
\end{definition}

\begin{proposition}[The post-$\mathbf K_2^{\mathrm{cur}}$ slot $\sigma_3$ admits only the transfer-positive or boundary-active status]
\label{prop:post-k2-slot-sigma3-admits-only-transfer-positive-or-boundary-active-status}
Assume the setting of Definition~\ref{def:transfer-positive-versus-boundary-active-status-at-sigma3}. Then exactly one of the following alternatives holds:
\begin{enumerate}[label=(\roman*)]
  \item $\sigma_3$ is transfer-positive, in which case the current line produces a valid three-step bounded certificate
  \[
    \mathbf K_3^{\mathrm{cand}}=(\mathcal C_0\subset \mathcal C_1\subset \mathcal C_2\subset \mathcal C_3;\,\sigma_1,\sigma_2,\sigma_3),
  \]
  with $\mathcal C_3$ the corridor obtained by adjoining $\sigma_3$ to $\mathcal C_2$; or
  \item $\sigma_3$ is boundary-active, in which case $\sigma_3$ itself is already the first genuine post-$\mathbf K_2^{\mathrm{cur}}$ boundary location.
\end{enumerate}
No third autonomous status for $\sigma_3$ is licensed on the present OP~1 extension line.
\end{proposition}

\begin{proof}
By Proposition~\ref{prop:first-honest-future-op1-extension-question-beyond-k2cur-localized-at-sigma3}, every honest future OP~1 extension question beyond $\mathbf K_2^{\mathrm{cur}}$ is localized at the single next leading admissible slot $\sigma_3$. Therefore the only remaining issue is whether $\sigma_3$ carries the same slotwise retained-discipline transfer certificate as the previously certified steps. If it does, then the bounded-$k$ certificate mechanism of Proposition~\ref{prop:bounded-k-step-certificates-certify-whole-finite-op1-run} extends the chain once more and yields the displayed three-step candidate certificate. If it does not, then by the same localization result $\sigma_3$ itself is already the first genuine post-$\mathbf K_2^{\mathrm{cur}}$ boundary location. Since the extension problem is completely localized at $\sigma_3$, these two alternatives exhaust the possibilities.
\end{proof}

\begin{corollary}[The next honest OP~1 decision beyond $\mathbf K_2^{\mathrm{cur}}$ is a two-channel decision at $\sigma_3$]
\label{cor:next-honest-op1-decision-beyond-k2cur-is-two-channel-decision-at-sigma3}
On the present OP~1 extension line, the first uncertified future decision beyond $\mathbf K_2^{\mathrm{cur}}$ reduces to the two-channel split
\[
  \text{transfer-positive at }\sigma_3
  \qquad\text{or}\qquad
  \text{boundary-active at }\sigma_3.
\]
In particular, the theorematic line and the Monte-Carlo/fusion-probe layer need not resurvey the already certified segment $\mathbf K_2^{\mathrm{cur}}$ in order to decide the next step.
\end{corollary}

\begin{proof}
Immediate from Proposition~\ref{prop:post-k2-slot-sigma3-admits-only-transfer-positive-or-boundary-active-status}.
\end{proof}

\begin{remark}[What the first transfer-positive/boundary-active split at $\sigma_3$ now buys]
\label{rem:what-first-transfer-positive-boundary-active-split-at-sigma3-now-buys}
This is the first point beyond the currently certified horizon where the next admissible slot is not merely isolated but also status-normalized. The benefit is that the next step no longer asks a vague continuation question. It asks only whether $\sigma_3$ is transfer-positive or already boundary-active. If it is transfer-positive, the line immediately upgrades to a concrete three-step certificate. If it is boundary-active, then the first genuine post-$\mathbf K_2^{\mathrm{cur}}$ boundary witness has already been found at the earliest possible place.
\end{remark}


\noindent\textbf{Reader cue.} Read the next block as the direct certification of the transfer-positive channel at $\sigma_3$. Its purpose is not yet to promote the extension line indefinitely, but to convert the already normalized positive branch into an explicit current three-step bounded certificate and to discharge the boundary-active channel at that slot itself.

\subsection*{19AZ$^{(6t)}$. First certification of $\mathbf K_3^{\mathrm{cur}}$ under transfer-positive status at $\sigma_3$}
\label{sec:first-certification-of-k3cur-under-transfer-positive-status-at-sigma3}
After 19AZ$^{(6s)}$, the next honest step is no longer to restate the two-channel split at $\sigma_3$, but to close its positive side. If the first later leading admissible slot beyond $\mathbf K_2^{\mathrm{cur}}$ is transfer-positive, then the same slotwise retained-discipline transfer mechanism that already produced the certified steps through $\sigma_2$ should certify one more step. The present block records that upgrade explicitly and removes the boundary-active alternative at $\sigma_3$ itself whenever the positive certificate is in force.

\begin{definition}[Present three-step bounded certificate $\mathbf K_3^{\mathrm{cur}}$]
\label{def:present-three-step-bounded-certificate-k3cur}
Assume the setting of Definition~\ref{def:transfer-positive-versus-boundary-active-status-at-sigma3}. If $\sigma_3$ is transfer-positive, define the \emph{present three-step bounded certificate}
\[
  \mathbf K_3^{\mathrm{cur}}=(\mathcal C_0\subset \mathcal C_1\subset \mathcal C_2\subset \mathcal C_3;\,\sigma_1,\sigma_2,\sigma_3),
\]
where $\mathcal C_3$ is obtained from $\mathcal C_2$ by adjoining the transfer-positive leading admissible slot $\sigma_3$ under the same slotwise retained-discipline transfer certificate used in the already certified steps.
\end{definition}

\begin{proposition}[Transfer-positive status at $\sigma_3$ certifies $\mathbf K_3^{\mathrm{cur}}$ and discharges the local boundary channel]
\label{prop:transfer-positive-status-at-sigma3-certifies-k3cur-and-discharges-local-boundary-channel}
Assume the setting of Definition~\ref{def:present-three-step-bounded-certificate-k3cur}. Then:
\begin{enumerate}[label=(\roman*)]
  \item the chain $\mathbf K_3^{\mathrm{cur}}$ is a valid bounded three-step OP~1 admissible-window certificate on the present extension line;
  \item $\sigma_3$ is not a genuine post-$\mathbf K_2^{\mathrm{cur}}$ boundary-witness location; and
  \item every honest future OP~1 extension failure, if any, must occur at a later leading admissible slot strictly beyond the corridor $\mathcal C_3$ certified by $\mathbf K_3^{\mathrm{cur}}$.
\end{enumerate}
\end{proposition}

\begin{proof}
By Definition~\ref{def:present-three-step-bounded-certificate-k3cur}, the slot $\sigma_3$ carries the same slotwise retained-discipline transfer certificate as the already certified steps through $\sigma_2$. Proposition~\ref{prop:bounded-k-step-certificates-certify-whole-finite-op1-run} therefore upgrades the chain by one step and certifies the displayed bounded three-step corridor, giving (i). Once that positive certification is in force, the alternative ``boundary-active at $\sigma_3$'' from Proposition~\ref{prop:post-k2-slot-sigma3-admits-only-transfer-positive-or-boundary-active-status} is excluded, proving (ii). Finally, after $\mathcal C_3$ has become part of the certified horizon, any later honest OP~1 failure must first appear strictly beyond $\mathcal C_3$, which is exactly (iii).
\end{proof}

\begin{corollary}[The positive channel at $\sigma_3$ upgrades the certified horizon from $\mathbf K_2^{\mathrm{cur}}$ to $\mathbf K_3^{\mathrm{cur}}$]
\label{cor:positive-channel-at-sigma3-upgrades-certified-horizon-from-k2cur-to-k3cur}
On the present OP~1 extension line, transfer-positive status at $\sigma_3$ does not merely produce a candidate continuation. It certifies the concrete bounded horizon $\mathbf K_3^{\mathrm{cur}}$ and pushes every honest future extension question to the first later leading admissible slot beyond $\mathcal C_3$.
\end{corollary}

\begin{proof}
Immediate from Proposition~\ref{prop:transfer-positive-status-at-sigma3-certifies-k3cur-and-discharges-local-boundary-channel}.
\end{proof}

\begin{remark}[What certifying $\mathbf K_3^{\mathrm{cur}}$ now buys]
\label{rem:what-certifying-k3cur-now-buys}
This is the first point at which the extension line no longer talks about a hypothetical third step but about a named current three-step horizon. The benefit is again twofold. The theorematic line now has a larger certified corridor and therefore a smaller honest future search zone. The Monte-Carlo/fusion-probe layer likewise no longer needs to test anything at $\sigma_3$ itself; its next bounded question starts only at the first later leading admissible slot beyond $\mathcal C_3$.
\end{remark}


\noindent\textbf{Reader cue.} Read the next block as the symmetric normalization of the negative channel at $\sigma_3$. Its purpose is not to reopen the already certified positive branch, but to record what the extension line is honestly allowed to say when $\sigma_3$ is boundary-active: namely, that the first explicit post-$\mathbf K_2^{\mathrm{cur}}$ extension-boundary witness has already been found there and that no later slot needs to be consulted before that witness is discharged.

\subsection*{19AZ$^{(6u)}$. First explicit normalization of the boundary-active channel at $\sigma_3$ as the earliest post-$\mathbf K_2^{\mathrm{cur}}$ extension witness}
\label{sec:first-explicit-normalization-of-boundary-active-channel-at-sigma3-as-earliest-post-k2cur-extension-witness}
After 19AZ$^{(6t)}$, the positive side at $\sigma_3$ has been normalized: if $\sigma_3$ is transfer-positive, then the line upgrades to the concrete bounded certificate $\mathbf K_3^{\mathrm{cur}}$. The negative side should now be stated with the same clarity. If $\sigma_3$ is boundary-active, then it is not merely a failed candidate slot. It is already the first explicit post-$\mathbf K_2^{\mathrm{cur}}$ extension witness, and every honest later extension question remains downstream of that earliest active boundary location until the witness at $\sigma_3$ has been discharged.

\begin{definition}[Current boundary-active witness datum at $\sigma_3$]
\label{def:current-boundary-active-witness-datum-at-sigma3}
Assume the setting of Definition~\ref{def:transfer-positive-versus-boundary-active-status-at-sigma3}. If $\sigma_3$ is boundary-active, define the \emph{current boundary-active witness datum}
\[
  \mathfrak B_3^{\mathrm{cur}}=(\mathcal C_2,\sigma_3),
\]
viewed as the earliest explicit post-$\mathbf K_2^{\mathrm{cur}}$ OP~1 extension-boundary witness on the present line.
\end{definition}

\begin{proposition}[Boundary-active status at $\sigma_3$ normalizes the earliest explicit post-$\mathbf K_2^{\mathrm{cur}}$ extension witness]
\label{prop:boundary-active-status-at-sigma3-normalizes-earliest-explicit-post-k2cur-extension-witness}
Assume the setting of Definition~\ref{def:current-boundary-active-witness-datum-at-sigma3}. Then:
\begin{enumerate}[label=(\roman*)]
  \item $\sigma_3$ is the first genuine post-$\mathbf K_2^{\mathrm{cur}}$ OP~1 extension-boundary location on the present line;
  \item no valid bounded three-step certificate $\mathbf K_3^{\mathrm{cur}}$ is licensed on the same line until the witness datum $\mathfrak B_3^{\mathrm{cur}}$ has been discharged; and
  \item every honest later OP~1 extension question remains downstream of $\sigma_3$ and need not be opened before the witness datum at $\sigma_3$ has been resolved.
\end{enumerate}
\end{proposition}

\begin{proof}
If $\sigma_3$ is boundary-active, then by Proposition~\ref{prop:post-k2-slot-sigma3-admits-only-transfer-positive-or-boundary-active-status} the positive alternative fails exactly at the first later leading admissible slot beyond $\mathbf K_2^{\mathrm{cur}}$. This gives (i). Since Proposition~\ref{prop:transfer-positive-status-at-sigma3-certifies-k3cur-and-discharges-local-boundary-channel} requires transfer-positive status in order to certify $\mathbf K_3^{\mathrm{cur}}$, the same line cannot simultaneously carry the certified three-step horizon while $\sigma_3$ remains boundary-active, proving (ii). Finally, because $\sigma_3$ is already the earliest active boundary location beyond the current certified horizon, no honest later extension question needs to be opened before that earliest witness has been addressed, which is (iii).
\end{proof}

\begin{corollary}[The negative channel at $\sigma_3$ is the earliest explicit witness beyond $\mathbf K_2^{\mathrm{cur}}$]
\label{cor:negative-channel-at-sigma3-is-earliest-explicit-witness-beyond-k2cur}
On the present OP~1 extension line, boundary-active status at $\sigma_3$ does not create a diffuse failure region. It produces the explicit earliest witness datum $\mathfrak B_3^{\mathrm{cur}}=(\mathcal C_2,\sigma_3)$ and blocks further honest extension promotion on that line until this witness has been discharged.
\end{corollary}

\begin{proof}
Immediate from Proposition~\ref{prop:boundary-active-status-at-sigma3-normalizes-earliest-explicit-post-k2cur-extension-witness}.
\end{proof}

\begin{remark}[What explicitly normalizing the boundary-active channel at $\sigma_3$ now buys]
\label{rem:what-explicitly-normalizing-boundary-active-channel-at-sigma3-now-buys}
This is the symmetric counterpart to certifying $\mathbf K_3^{\mathrm{cur}}$ on the positive side. The benefit is that the negative side is no longer left as an unnamed failure possibility. It is packaged as the earliest explicit post-$\mathbf K_2^{\mathrm{cur}}$ witness datum, so that the theorematic line, the routing layer, and the Monte-Carlo/fusion-probe layer all know where the next honest obstruction sits if the transfer-positive channel does not carry. In particular, no later slot needs to be searched before $\sigma_3$ itself has been resolved.
\end{remark}


\noindent\textbf{Reader cue.} Read the next block as the compact decision-normal form for the first slot beyond the current certified horizon. Its purpose is not to add a new extension theorem, but to compress the already normalized alternatives at $\sigma_3$ into one finite decision object that the theorematic line, the routing layer, and the probe layer can all use without reopening the full case split.

\subsection*{19AZ$^{(6v)}$. First decision-normal form at $\sigma_3$ beyond the current certified horizon}
\label{sec:first-decision-normal-form-at-sigma3-beyond-current-certified-horizon}
After 19AZ$^{(6t)}$ and 19AZ$^{(6u)}$, the first slot beyond the current certified horizon has been normalized on both sides. On the positive side, transfer-positive status at $\sigma_3$ certifies the concrete bounded three-step horizon $\mathbf K_3^{\mathrm{cur}}$. On the negative side, boundary-active status at $\sigma_3$ produces the earliest explicit post-$\mathbf K_2^{\mathrm{cur}}$ witness datum $\mathfrak B_3^{\mathrm{cur}}=(\mathcal C_2,\sigma_3)$. The next honest step is therefore to package these two outcomes into one compact decision-normal form for the first later leading slot beyond the presently certified horizon.

\begin{definition}[Current horizon decision datum at $\sigma_3$]
\label{def:current-horizon-decision-datum-at-sigma3}
Assume the setting of Definition~\ref{def:transfer-positive-versus-boundary-active-status-at-sigma3}. Define the \emph{current horizon decision datum}
\[
  \mathfrak D_3^{\mathrm{cur}}\in\{\mathbf K_3^{\mathrm{cur}},\mathfrak B_3^{\mathrm{cur}}\}
\]
by the rule
\[
  \mathfrak D_3^{\mathrm{cur}}=
  \begin{cases}
    \mathbf K_3^{\mathrm{cur}}, & \text{if }\sigma_3\text{ is transfer-positive},\\[2mm]
    \mathfrak B_3^{\mathrm{cur}}, & \text{if }\sigma_3\text{ is boundary-active}.
  \end{cases}
\]
\end{definition}

\begin{proposition}[The first slot beyond $\mathbf K_2^{\mathrm{cur}}$ admits a finite decision-normal form]
\label{prop:first-slot-beyond-k2cur-admits-finite-decision-normal-form}
Assume the setting of Definition~\ref{def:current-horizon-decision-datum-at-sigma3}. Then:
\begin{enumerate}[label=(\roman*)]
  \item the datum $\mathfrak D_3^{\mathrm{cur}}$ is well-defined and exhaustive for the first later leading admissible slot beyond $\mathbf K_2^{\mathrm{cur}}$;
  \item if $\mathfrak D_3^{\mathrm{cur}}=\mathbf K_3^{\mathrm{cur}}$, then the certified horizon has advanced through $\mathcal C_3$ and every honest later extension question begins strictly beyond that corridor; and
  \item if $\mathfrak D_3^{\mathrm{cur}}=\mathfrak B_3^{\mathrm{cur}}$, then $\sigma_3$ is already the earliest explicit post-$\mathbf K_2^{\mathrm{cur}}$ boundary location and no later slot needs to be opened before that witness datum has been discharged.
\end{enumerate}
\end{proposition}

\begin{proof}
By Proposition~\ref{prop:post-k2-slot-sigma3-admits-only-transfer-positive-or-boundary-active-status}, the slot $\sigma_3$ admits only the two statuses transfer-positive and boundary-active. Definition~\ref{def:current-horizon-decision-datum-at-sigma3} therefore assigns exactly one of the two normalized outputs $\mathbf K_3^{\mathrm{cur}}$ and $\mathfrak B_3^{\mathrm{cur}}$, proving (i). If the positive output occurs, Proposition~\ref{prop:transfer-positive-status-at-sigma3-certifies-k3cur-and-discharges-local-boundary-channel} gives the certified three-step horizon and pushes every honest later extension question beyond $\mathcal C_3$, which is (ii). If the negative output occurs, Proposition~\ref{prop:boundary-active-status-at-sigma3-normalizes-earliest-explicit-post-k2cur-extension-witness} gives the earliest explicit witness datum at $\sigma_3$ and forbids opening later slots first, which is (iii).
\end{proof}

\begin{corollary}[The first slot beyond the current certified horizon is now decision-normalized]
\label{cor:first-slot-beyond-current-certified-horizon-is-now-decision-normalized}
On the present OP~1 extension line, the first later leading admissible slot beyond $\mathbf K_2^{\mathrm{cur}}$ no longer carries an uncompressed continuation question. It is represented by the finite decision datum $\mathfrak D_3^{\mathrm{cur}}\in\{\mathbf K_3^{\mathrm{cur}},\mathfrak B_3^{\mathrm{cur}}\}$, which either upgrades the current certified horizon or identifies the earliest explicit boundary witness.
\end{corollary}

\begin{proof}
Immediate from Proposition~\ref{prop:first-slot-beyond-k2cur-admits-finite-decision-normal-form}.
\end{proof}

\begin{remark}[What the decision-normal form at $\sigma_3$ now buys]
\label{rem:what-decision-normal-form-at-sigma3-now-buys}
This is the first point beyond the currently certified horizon where the extension line can speak in one compact finite decision object instead of in a longer two-branch case split. The theorematic line gains a cleaner handoff to the next extension stage; the routing layer gains a single certified output to follow; and the Monte-Carlo/fusion-probe layer gains a sharply bounded question: either the line has already promoted to $\mathbf K_3^{\mathrm{cur}}$, or the earliest active obstruction has already been localized at $\sigma_3$ itself.
\end{remark}




\noindent\textbf{Reader cue.} Read the next block as the compact QAM/decoder output language for the first extension decision beyond the current certified horizon. Its purpose is not to add a stronger extension theorem, but to package the already normalized alternatives \(\mathbf K_3^{\mathrm{cur}}\) versus \(\mathfrak B_3^{\mathrm{cur}}\) into one short token language that later theorematic, routing, and probe layers can cite without reopening the full decision datum.

\subsection*{19AZ$^{(6w)}$. First QAM/decoder output tokens for extension tasks beyond the current certified horizon}
\label{sec:first-qam-decoder-output-tokens-for-extension-tasks-beyond-current-certified-horizon}
After 19AZ$^{(6v)}$, the first slot beyond the current certified horizon is represented by the finite decision datum
\[
  \mathfrak D_3^{\mathrm{cur}}\in\{\mathbf K_3^{\mathrm{cur}},\mathfrak B_3^{\mathrm{cur}}\}.
\]
The next honest compression step is therefore to translate this extension decision into a short output language that can be consumed uniformly by the QAM layer, the route selector, and the probe/decoder layer without reopening the full theorematic split.

\begin{definition}[Extension-output token language at $\sigma_3$]
\label{def:extension-output-token-language-at-sigma3}
Define the finite token set
\[
  \mathcal T^{\mathrm{ext}}_3:=\{\mathsf{EK}_3,\mathsf{EB}_3\},
\]
where the symbols mean:
\begin{align*}
  \mathsf{EK}_3 &:= \text{``the current horizon extends through }\mathcal C_3\text{''},\\
  \mathsf{EB}_3 &:= \text{``the earliest honest post-}\mathbf K_2^{\mathrm{cur}}\text{ boundary sits at }\sigma_3\text{''}.
\end{align*}
Define the associated token map by
\[
  \mathsf{Out}^{\mathrm{tok}}_{\mathrm{ext},3}(\mathfrak D_3^{\mathrm{cur}})=
  \begin{cases}
    \mathsf{EK}_3, & \text{if }\mathfrak D_3^{\mathrm{cur}}=\mathbf K_3^{\mathrm{cur}},\\[2mm]
    \mathsf{EB}_3, & \text{if }\mathfrak D_3^{\mathrm{cur}}=\mathfrak B_3^{\mathrm{cur}}.
  \end{cases}
\]
\end{definition}

\begin{proposition}[The first extension decision beyond $\mathbf K_2^{\mathrm{cur}}$ admits a faithful token output]
\label{prop:first-extension-decision-beyond-k2cur-admits-faithful-token-output}
Assume the setting of Definition~\ref{def:extension-output-token-language-at-sigma3}. Then:
\begin{enumerate}[label=(\roman*)]
  \item the map \(\mathsf{Out}^{\mathrm{tok}}_{\mathrm{ext},3}\) is well-defined on the current decision datum \(\mathfrak D_3^{\mathrm{cur}}\);
  \item \(\mathsf{Out}^{\mathrm{tok}}_{\mathrm{ext},3}(\mathfrak D_3^{\mathrm{cur}})=\mathsf{EK}_3\) if and only if the line has already promoted to the certified bounded horizon \(\mathbf K_3^{\mathrm{cur}}\); and
  \item \(\mathsf{Out}^{\mathrm{tok}}_{\mathrm{ext},3}(\mathfrak D_3^{\mathrm{cur}})=\mathsf{EB}_3\) if and only if \(\sigma_3\) is already the earliest explicit boundary witness beyond \(\mathbf K_2^{\mathrm{cur}}\).
\end{enumerate}
In particular, the token output is faithful to the full current extension decision.
\end{proposition}

\begin{proof}
By Definition~\ref{def:current-horizon-decision-datum-at-sigma3}, the decision datum \(\mathfrak D_3^{\mathrm{cur}}\) takes exactly one of the two values \(\mathbf K_3^{\mathrm{cur}}\) and \(\mathfrak B_3^{\mathrm{cur}}\), so the token map is well-defined, proving (i). If the first value occurs, then Proposition~\ref{prop:transfer-positive-status-at-sigma3-certifies-k3cur-and-discharges-local-boundary-channel} shows that the certified horizon has advanced through \(\mathcal C_3\), hence the output token is exactly \(\mathsf{EK}_3\), proving (ii). If the second value occurs, then Proposition~\ref{prop:boundary-active-status-at-sigma3-normalizes-earliest-explicit-post-k2cur-extension-witness} shows that \(\sigma_3\) is the earliest explicit post-\(\mathbf K_2^{\mathrm{cur}}\) boundary location, hence the output token is exactly \(\mathsf{EB}_3\), proving (iii). Since the two tokens are distinct and attached to the two distinct normalized outcomes, the output is faithful to the full decision datum.
\end{proof}

\begin{corollary}[Extension tasks at the first post-$\mathbf K_2^{\mathrm{cur}}$ slot are now token-normalized]
\label{cor:extension-tasks-at-first-post-k2cur-slot-are-now-token-normalized}
On the present OP~1 extension line, the first extension task beyond the certified horizon \(\mathbf K_2^{\mathrm{cur}}\) can now be recorded by the compact token
\[
  \mathsf{Out}^{\mathrm{tok}}_{\mathrm{ext},3}(\mathfrak D_3^{\mathrm{cur}})\in\{\mathsf{EK}_3,\mathsf{EB}_3\}.
\]
Thus later routing or probe logic no longer needs the full branch split at \(\sigma_3\); it is enough to ask whether the current token is \(\mathsf{EK}_3\) or \(\mathsf{EB}_3\).
\end{corollary}

\begin{proof}
Immediate from Proposition~\ref{prop:first-extension-decision-beyond-k2cur-admits-faithful-token-output}.
\end{proof}

\begin{remark}[What the first extension-output tokens now buy]
\label{rem:what-first-extension-output-tokens-now-buy}
The gain is small but real. The present line no longer needs to treat the first post-\(\mathbf K_2^{\mathrm{cur}}\) extension task as a longer prose decision between ``certify the next horizon'' and ``localize the earliest new witness''. The theorematic layer, the routing layer, and the Monte-Carlo/probe layer may all read the same compact output token \(\mathsf{EK}_3\) or \(\mathsf{EB}_3\). In particular, this is the first point where the extension line itself begins to speak the same kind of compact QAM/decoder language that had previously been reserved for reopening tasks inside the retained OP kernels.
\end{remark}



\noindent\textbf{Reader cue.} Read the next block as the reusable proof-compression template abstracted from the first post-$\mathbf K_2^{\mathrm{cur}}$ decision. Its purpose is not to claim a new global extension theorem, but to isolate the exact finite-decision pattern that any later first-slot extension stage must satisfy once its local positive-versus-boundary dichotomy has been established.

\subsection*{19AZ$^{(6x)}$. Reusable next-slot decision schema beyond a certified horizon}
\label{sec:reusable-next-slot-decision-schema-beyond-certified-horizon}
The blocks 19AZ$^{(6v)}$ and 19AZ$^{(6w)}$ show that, once the first later leading admissible slot beyond a certified current horizon has been normalized into the two statuses \emph{transfer-positive} versus \emph{boundary-active}, the continuation question collapses to one finite decision datum and one faithful output token. The point of the present block is to isolate that pattern as a reusable local schema.

\begin{definition}[Reusable next-slot decision schema]
\label{def:reusable-next-slot-decision-schema}
Fix a current certified horizon $\mathbf K_m^{\mathrm{cur}}$, and let $\sigma_{m+1}$ denote the first later leading admissible slot beyond it. Assume that at this slot the local theorematic line has already established an exhaustive dichotomy:
\[
\text{$\sigma_{m+1}$ is either transfer-positive or boundary-active.}
\]
Assume further that:
\begin{itemize}[leftmargin=2em]
  \item transfer-positive status at $\sigma_{m+1}$ certifies a concrete upgraded horizon $\mathbf K_{m+1}^{\mathrm{cur}}$, and
  \item boundary-active status at $\sigma_{m+1}$ produces an earliest explicit witness datum $\mathfrak B_{m+1}^{\mathrm{cur}}$.
\end{itemize}
Then define the associated \emph{next-slot decision datum}
\[
  \mathfrak D_{m+1}^{\mathrm{cur}}
  \in
  \{\mathbf K_{m+1}^{\mathrm{cur}},\mathfrak B_{m+1}^{\mathrm{cur}}\}
\]
by
\[
  \mathfrak D_{m+1}^{\mathrm{cur}}
  =
  \begin{cases}
    \mathbf K_{m+1}^{\mathrm{cur}}, & \text{if }\sigma_{m+1}\text{ is transfer-positive},\\[2mm]
    \mathfrak B_{m+1}^{\mathrm{cur}}, & \text{if }\sigma_{m+1}\text{ is boundary-active}.
  \end{cases}
\]
Define moreover the associated finite output-token set
\[
  \mathcal T^{\mathrm{ext}}_{m+1}
  :=
  \{\mathsf{EK}_{m+1},\mathsf{EB}_{m+1}\},
\]
where
\begin{align*}
  \mathsf{EK}_{m+1} &:= \text{``the current horizon extends through }\mathcal C_{m+1}\text{''},\\
  \mathsf{EB}_{m+1} &:= \text{``the earliest honest boundary beyond }\mathbf K_m^{\mathrm{cur}}\text{ sits at }\sigma_{m+1}\text{''}.
\end{align*}
Finally define the token output map
\[
  \mathsf{Out}^{\mathrm{tok}}_{\mathrm{ext},m+1}(\mathfrak D_{m+1}^{\mathrm{cur}})
  =
  \begin{cases}
    \mathsf{EK}_{m+1}, & \text{if }\mathfrak D_{m+1}^{\mathrm{cur}}=\mathbf K_{m+1}^{\mathrm{cur}},\\[2mm]
    \mathsf{EB}_{m+1}, & \text{if }\mathfrak D_{m+1}^{\mathrm{cur}}=\mathfrak B_{m+1}^{\mathrm{cur}}.
  \end{cases}
\]
\end{definition}

\begin{proposition}[Any normalized first-slot extension stage collapses to one finite decision datum and one faithful token]
\label{prop:any-normalized-first-slot-extension-stage-collapses-to-one-finite-decision-datum-and-one-faithful-token}
Assume the setting of Definition~\ref{def:reusable-next-slot-decision-schema}. Then:
\begin{enumerate}[label=(\roman*)]
  \item the datum $\mathfrak D_{m+1}^{\mathrm{cur}}$ is well-defined and exhaustive for the first later leading admissible slot beyond $\mathbf K_m^{\mathrm{cur}}$;
  \item the token map $\mathsf{Out}^{\mathrm{tok}}_{\mathrm{ext},m+1}$ is well-defined and faithful to that decision datum; and
  \item every later routing, decoder, or probe layer may consume the compact token $\mathsf{EK}_{m+1}$ or $\mathsf{EB}_{m+1}$ without reopening the full local branch split at $\sigma_{m+1}$.
\end{enumerate}
\end{proposition}

\begin{proof}
By assumption, the local line at $\sigma_{m+1}$ admits exactly the two statuses transfer-positive and boundary-active, and each status produces exactly one normalized output, namely $\mathbf K_{m+1}^{\mathrm{cur}}$ or $\mathfrak B_{m+1}^{\mathrm{cur}}$. This proves (i). The token map is then defined on exactly these two outputs and assigns distinct tokens to them, so it is well-defined and faithful, proving (ii). Since the token remembers which of the two normalized outcomes occurred, any later routing, decoder, or probe layer may consume the token alone unless it explicitly chooses to reopen the local theorematic proof, proving (iii).
\end{proof}

\begin{corollary}[The $\sigma_3$ decision package is the first realized instance of the reusable extension schema]
\label{cor:sigma3-decision-package-is-the-first-realized-instance-of-reusable-extension-schema}
The package
\[
  \mathfrak D_3^{\mathrm{cur}}
  \in
  \{\mathbf K_3^{\mathrm{cur}},\mathfrak B_3^{\mathrm{cur}}\},
  \qquad
  \mathsf{Out}^{\mathrm{tok}}_{\mathrm{ext},3}
  \in
  \{\mathsf{EK}_3,\mathsf{EB}_3\},
\]
constructed in 19AZ$^{(6v)}$--19AZ$^{(6w)}$, is not merely a local convenience. It is the first realized instance of the general next-slot decision schema for extension tasks beyond a certified current horizon.
\end{corollary}

\begin{proof}
Immediate from Definition~\ref{def:reusable-next-slot-decision-schema} and Proposition~\ref{prop:any-normalized-first-slot-extension-stage-collapses-to-one-finite-decision-datum-and-one-faithful-token} by taking $m=2$.
\end{proof}

\begin{remark}[What this schema now buys]
\label{rem:what-this-schema-now-buys}
This is the point where the present OP~1 extension line begins to export a genuine reusable tool rather than only a local closure fact. Future first-slot extension stages no longer need to reinvent the packaging logic: once the local positive-versus-boundary dichotomy is proved, the decision datum and faithful output token come for free. In particular, the current $\sigma_3$ block is now readable as the prototype of a general extension-normalization interface.
\end{remark}


\noindent\textbf{Reader cue.} Read the next block as the reusable action map generated by the extension-token language. Its purpose is not to add a stronger horizon theorem, but to force every later layer to take one of only two honest next actions: either advance the certified horizon, or reopen work exactly at the earliest normalized boundary site.

\subsection*{19AZ$^{(6y)}$. Reusable action/transition map driven by next-slot extension tokens}
\label{sec:reusable-action-transition-map-driven-by-next-slot-extension-tokens}
After 19AZ$^{(6x)}$, the line has not only a reusable next-slot decision datum but also a reusable token language. The next honest compression step is therefore to attach to that token language the corresponding next action. The guiding principle is simple: a token should not merely summarize the local result, but should also determine which later theorematic, routing, decoder, or probe move is licensed and which moves are forbidden.

\begin{definition}[Reusable extension-action map]
\label{def:reusable-extension-action-map}
Assume the setting of Definition~\ref{def:reusable-next-slot-decision-schema}. Define the finite extension-action set
\[
  \mathcal A^{\mathrm{ext}}_{m+1}
  :=
  \{\mathsf{ADV}_{m+1},\mathsf{OPENB}_{m+1}\},
\]
where the symbols mean:
\begin{align*}
  \mathsf{ADV}_{m+1} &:= \text{``advance the current certified horizon through }\mathcal C_{m+1}\text{ and shift the first unresolved extension question strictly beyond it''},\\
  \mathsf{OPENB}_{m+1} &:= \text{``open the earliest normalized boundary kernel exactly at }\sigma_{m+1}\text{ and do not skip to later slots first''}.
\end{align*}
Define the associated action map on extension tokens by
\[
  \mathsf{Act}^{\mathrm{tok}}_{\mathrm{ext},m+1}(\tau)
  =
  \begin{cases}
    \mathsf{ADV}_{m+1}, & \text{if }\tau=\mathsf{EK}_{m+1},\\[2mm]
    \mathsf{OPENB}_{m+1}, & \text{if }\tau=\mathsf{EB}_{m+1}.
  \end{cases}
\]
Equivalently, one may read the composite transition map
\[
  \mathsf{Act}^{\mathrm{tok}}_{\mathrm{ext},m+1}\circ \mathsf{Out}^{\mathrm{tok}}_{\mathrm{ext},m+1}
\]
as the minimal action normal form attached to the first later leading admissible slot beyond \(\mathbf K_m^{\mathrm{cur}}\).
\end{definition}

\begin{proposition}[The extension-token language determines the next honest action]
\label{prop:extension-token-language-determines-next-honest-action}
Assume the setting of Definition~\ref{def:reusable-extension-action-map}. Then:
\begin{enumerate}[label=(\roman*)]
  \item the map \(\mathsf{Act}^{\mathrm{tok}}_{\mathrm{ext},m+1}\) is well-defined on \(\mathcal T^{\mathrm{ext}}_{m+1}\);
  \item if the output action is \(\mathsf{ADV}_{m+1}\), then the current line is licensed to replace \(\mathbf K_m^{\mathrm{cur}}\) by \(\mathbf K_{m+1}^{\mathrm{cur}}\) and to relocate the first unresolved extension question strictly beyond \(\mathcal C_{m+1}\);
  \item if the output action is \(\mathsf{OPENB}_{m+1}\), then every honest later theorematic, routing, decoder, or probe step must remain localized at the earliest normalized boundary site \(\sigma_{m+1}\) until that local boundary kernel has been processed.
\end{enumerate}
\end{proposition}

\begin{proof}
The token language from Definition~\ref{def:reusable-next-slot-decision-schema} has exactly the two values \(\mathsf{EK}_{m+1}\) and \(\mathsf{EB}_{m+1}\), so the action map is well-defined, proving (i). If the first token occurs, then by construction it records that the current horizon has already advanced through \(\mathcal C_{m+1}\), so the only honest next action is to promote the certified horizon and move the unresolved extension question strictly beyond that corridor, proving (ii). If the second token occurs, then by construction it records that the earliest honest obstruction has already been normalized at \(\sigma_{m+1}\); therefore any later layer that skips ahead before processing that site would violate the earliest-boundary discipline, proving (iii).
\end{proof}

\begin{corollary}[The \(\sigma_3\) token package now induces a certified next-action split]
\label{cor:sigma3-token-package-now-induces-certified-next-action-split}
For the realized package from 19AZ$^{(6v)}$--19AZ$^{(6x)}$, the composite map
\[
  \mathsf{Act}^{\mathrm{tok}}_{\mathrm{ext},3}\circ \mathsf{Out}^{\mathrm{tok}}_{\mathrm{ext},3}
\]
returns exactly one of the two certified next actions:
\[
  \mathsf{ADV}_3
  \qquad\text{or}\qquad
  \mathsf{OPENB}_3.
\]
Thus the present line is now action-normalized at the first slot beyond \(\mathbf K_2^{\mathrm{cur}}\): either it advances through \(\mathcal C_3\), or it reopens work exactly at the earliest explicit boundary site \(\sigma_3\).
\end{corollary}

\begin{proof}
Immediate from Proposition~\ref{prop:extension-token-language-determines-next-honest-action} by taking \(m=2\).
\end{proof}

\begin{remark}[What this action map now buys, and what can be salvaged into later layers]
\label{rem:what-this-action-map-now-buys-and-what-can-be-salvaged-into-later-layers}
This is the point where the extension-tool layer stops being a passive encoding device and becomes a genuine routing interface. Later consumers no longer need to ask the full local question again; they need only read whether the current action is \(\mathsf{ADV}_{m+1}\) or \(\mathsf{OPENB}_{m+1}\). That is precisely the salvageable discipline for the later optimization ideas: Monte-Carlo/fusion probes, decoder packages, and any later Heisenberg-matrix, Millennium-matrix, crystal-equation, or field-equation diagnostics should attach below this action interface as consumers of a certified local action, not as mechanisms that reopen or overwrite the theorematic decision itself.
\end{remark}


\noindent\textbf{Reader cue.} Read the next block as the iterative upgrade of the local action map. Its purpose is not to claim an already completed global induction, but to isolate the exact protocol by which the active line may repeatedly consume one certified next-slot token at a time, advancing when the token says ``extend'' and stopping honestly when the token says ``boundary''.

\subsection*{19AZ$^{(6z)}$. Iterated extension protocol generated by certified local action tokens}
\label{sec:iterated-extension-protocol-generated-by-certified-local-action-tokens}
After 19AZ$^{(6y)}$, the line has a local action interface for the first later leading admissible slot beyond a certified current horizon. The next honest compression step is to package that interface into an iterative protocol. The point is simple: once one local stage has been normalized into a token and then into an action, the same logic can be reused at the next stage provided the new local dichotomy is again theorematically certified. Thus the extension line should be readable not as a one-off move at $\sigma_3$, but as a repeatable one-slot protocol with a built-in stop rule.

\begin{definition}[Iterated extension protocol with earliest-boundary stop rule]
\label{def:iterated-extension-protocol-with-earliest-boundary-stop-rule}
Fix a certified current horizon $\mathbf K_m^{\mathrm{cur}}$. Suppose that for each stage $r\geq m+1$, whenever the line reaches a certified horizon $\mathbf K_{r-1}^{\mathrm{cur}}$, the first later leading admissible slot $\sigma_r$ admits a certified local dichotomy of the kind required in Definition~\ref{def:reusable-next-slot-decision-schema}. Then define the associated stage-$r$ token
\[
  \tau_r \in \mathcal T^{\mathrm{ext}}_r = \{\mathsf{EK}_r,\mathsf{EB}_r\}
\]
and stage-$r$ action
\[
  a_r := \mathsf{Act}^{\mathrm{tok}}_{\mathrm{ext},r}(\tau_r)
  \in \mathcal A^{\mathrm{ext}}_r = \{\mathsf{ADV}_r,\mathsf{OPENB}_r\}.
\]
The resulting \emph{iterated extension protocol} is the finite or infinite stage-by-stage procedure that obeys the following rule at each step:
\begin{enumerate}[label=(\roman*)]
  \item if $a_r=\mathsf{ADV}_r$, replace the current certified horizon $\mathbf K_{r-1}^{\mathrm{cur}}$ by $\mathbf K_r^{\mathrm{cur}}$ and continue to the next stage $r+1$;
  \item if $a_r=\mathsf{OPENB}_r$, stop the forward extension protocol and localize all further honest work at the earliest normalized boundary kernel sitting at $\sigma_r$.
\end{enumerate}
The first index at which the protocol stops, if such an index exists, is called the \emph{earliest protocol boundary stage} and is denoted by
\[
  r_* := \min\{r\ge m+1 : a_r=\mathsf{OPENB}_r\}.
\]
When no such stage exists in the range under consideration, the protocol is said to be \emph{advance-persistent} on that range.
\end{definition}

\begin{proposition}[The iterated extension protocol is well-defined and locally non-skipping]
\label{prop:iterated-extension-protocol-is-well-defined-and-locally-non-skipping}
Assume the setting of Definition~\ref{def:iterated-extension-protocol-with-earliest-boundary-stop-rule}. Then:
\begin{enumerate}[label=(\roman*)]
  \item at each certified stage $r$, exactly one action $a_r\in\{\mathsf{ADV}_r,\mathsf{OPENB}_r\}$ is produced;
  \item if the protocol is stopped at stage $r_*$, then $\sigma_{r_*}$ is the earliest stage in the protocol at which the line is licensed to open a normalized boundary kernel rather than advance the certified horizon;
  \item no honest theorematic, routing, decoder, probe, or diagnostic layer may skip past a stage with action $\mathsf{OPENB}_r$ and continue the forward extension protocol as if an $\mathsf{ADV}_r$ action had occurred.
\end{enumerate}
\end{proposition}

\begin{proof}
For each certified stage $r$, Definition~\ref{def:reusable-next-slot-decision-schema} yields exactly one token in $\{\mathsf{EK}_r,\mathsf{EB}_r\}$, and Definition~\ref{def:reusable-extension-action-map} sends that token to exactly one action in $\{\mathsf{ADV}_r,\mathsf{OPENB}_r\}$, proving (i). If the protocol stops at stage $r_*$, then by definition $r_*$ is the least stage whose action is $\mathsf{OPENB}_{r_*}$, so all earlier stages were genuine advance stages and $\sigma_{r_*}$ is the earliest protocol boundary stage, proving (ii). Statement (iii) is then forced by the meaning of $\mathsf{OPENB}_r$: it records that the earliest honest unresolved obstruction has already been normalized at stage $r$, so any layer that skipped past it would violate the non-skipping boundary discipline built into 19AZ$^{(6y)}$.
\end{proof}

\begin{corollary}[The present $\sigma_3$ package is the first step of an iterated extension protocol]
\label{cor:present-sigma3-package-is-first-step-of-iterated-extension-protocol}
In the current active line, the package from 19AZ$^{(6v)}$--19AZ$^{(6y)}$ is not only a local closure package at the first slot beyond $\mathbf K_2^{\mathrm{cur}}$. It is the first realized stage of the iterated extension protocol. Concretely, the line now carries a certified first-stage token
\[
  \tau_3 \in \{\mathsf{EK}_3,\mathsf{EB}_3\}
\]
and the associated first-stage action
\[
  a_3 \in \{\mathsf{ADV}_3,\mathsf{OPENB}_3\},
\]
so every later extension attempt must either continue from the promoted horizon $\mathbf K_3^{\mathrm{cur}}$ or remain localized at the earliest normalized boundary kernel at $\sigma_3$.
\end{corollary}

\begin{proof}
Immediate from Corollary~\ref{cor:sigma3-token-package-now-induces-certified-next-action-split} and Proposition~\ref{prop:iterated-extension-protocol-is-well-defined-and-locally-non-skipping}.
\end{proof}

\begin{remark}[What can be salvaged from later optimization ideas without corrupting the proof line]
\label{rem:what-can-be-salvaged-from-later-optimization-ideas-without-corrupting-proof-line}
This protocol is also the right place to salvage the useful part of the later optimization discussion. Once the active line emits a certified action stream $(a_r)$, later consumers may compress, hash, decode, profile, or statistically probe that stream --- for example through Monte-Carlo/fusion diagnostics, decoder packages, Markov-style routing summaries, or matrix-side readout layers --- provided they remain downstream consumers of certified stage actions. What they may \emph{not} do is replace a missing local dichotomy proof by a heuristic token guess, or overwrite an $\mathsf{OPENB}_r$ stop signal by pretending that a later advance has already been theorematically licensed.
\end{remark}


\noindent\textbf{Reader cue.} Read the next block as the finite word-level compression of the iterated extension protocol. Its purpose is not to replace the stagewise theorematic decisions, but to encode any already certified finite protocol segment into a short QAM/decoder-readable extension word whose shape itself records whether the segment remained purely advancing or terminated at its first honest boundary stage.

\subsection*{19AZ$^{(6aa)}$. Finite QAM/decoder word form for certified extension-action segments}
\label{sec:finite-qam-decoder-word-form-for-certified-extension-action-segments}
After 19AZ$^{(6z)}$, the active line has a stagewise protocol, but later consumers may still want a shorter object than the full indexed action stream. The next honest compression step is therefore to pass from certified stage actions to a finite extension word. The compression must remain faithful to the non-skipping boundary discipline: once an earliest boundary action appears, the forward word must terminate there rather than pretending to continue past that point.

\begin{definition}[Finite extension word alphabet and encoder]
\label{def:finite-extension-word-alphabet-and-encoder}
Fix a starting certified horizon \(\mathbf K_m^{\mathrm{cur}}\). Introduce the two-letter extension alphabet
\[
  \Sigma_{\mathrm{ext}} := \{\mathsf A,\mathsf B\},
\]
read as
\[
  \mathsf A = \text{``advance the certified horizon by one certified stage''},
  \qquad
  \mathsf B = \text{``open the earliest certified boundary kernel and stop''}.
\]
For each stage \(r\ge m+1\), define the stage encoder
\[
  \mathsf{enc}_r:\mathcal A^{\mathrm{ext}}_r\to\Sigma_{\mathrm{ext}}
\]
by
\[
  \mathsf{enc}_r(\mathsf{ADV}_r)=\mathsf A,
  \qquad
  \mathsf{enc}_r(\mathsf{OPENB}_r)=\mathsf B.
\]
If a certified protocol segment runs from stage \(m+1\) through stage \(N\), define its associated \emph{finite extension word}
\[
  \mathbf w_{m+1:N}
  :=
  \mathsf{enc}_{m+1}(a_{m+1})\,\mathsf{enc}_{m+2}(a_{m+2})\cdots\mathsf{enc}_N(a_N)
  \in \Sigma_{\mathrm{ext}}^{*},
\]
with the convention that if some stage \(r_*\le N\) satisfies \(a_{r_*}=\mathsf{OPENB}_{r_*}\), then the word is truncated at \(r_*\):
\[
  \mathbf w_{m+1:N}
  :=
  \mathsf{enc}_{m+1}(a_{m+1})\cdots\mathsf{enc}_{r_*}(a_{r_*}),
\]
because the protocol is not licensed to advance past the earliest boundary stage.
\end{definition}

\begin{proposition}[Certified extension words have exactly the form \(\mathsf A^k\) or \(\mathsf A^k\mathsf B\)]
\label{prop:certified-extension-words-have-exactly-form-Ak-or-AkB}
Assume the setting of Definition~\ref{def:iterated-extension-protocol-with-earliest-boundary-stop-rule}. Then every finite certified extension word produced by Definition~\ref{def:finite-extension-word-alphabet-and-encoder} has exactly one of the following forms:
\[
  \mathbf w = \mathsf A^k
  \qquad\text{or}\qquad
  \mathbf w = \mathsf A^k\mathsf B
\]
for some integer \(k\ge 0\). In particular:
\begin{enumerate}[label=(\roman*)]
  \item no certified extension word contains two occurrences of \(\mathsf B\);
  \item no certified extension word contains \(\mathsf B\) followed by any later symbol;
  \item the word alone already records whether the observed certified segment is advance-persistent on its range or has terminated at its earliest honest boundary stage.
\end{enumerate}
\end{proposition}

\begin{proof}
By Proposition~\ref{prop:iterated-extension-protocol-is-well-defined-and-locally-non-skipping}, each certified stage contributes exactly one action, either \(\mathsf{ADV}_r\) or \(\mathsf{OPENB}_r\). Under the encoder of Definition~\ref{def:finite-extension-word-alphabet-and-encoder}, these become \(\mathsf A\) and \(\mathsf B\). If no stage in the observed range opens a boundary kernel, then every action is an advance action and the resulting word is \(\mathsf A^k\) for the appropriate length \(k\). If a boundary stage occurs, let \(r_*\) be the earliest one. The protocol stops at \(r_*\), so the word is truncated there and has the form \(\mathsf A^k\mathsf B\), where \(k\) counts the earlier advance stages. This proves the displayed classification, and clauses (i)--(iii) are immediate consequences of that classification.
\end{proof}

\begin{corollary}[The extension protocol defines a regular terminal-word language]
\label{cor:extension-protocol-defines-regular-terminal-word-language}
On every finite observation range, the certified extension words belong to the terminal language
\[
  \mathcal L_{\mathrm{ext}}^{\mathrm{term}}
  =
  \{\mathsf A^k : k\ge 0\}
  \cup
  \{\mathsf A^k\mathsf B : k\ge 0\}.
\]
Hence every later QAM, decoder, hashing, routing-summary, or Monte-Carlo/profile layer that consumes certified extension words must respect this language and may not treat words of the forbidden forms
\[
  \mathsf B\mathsf A,\qquad \mathsf A\mathsf B\mathsf A,\qquad \mathsf B\mathsf B,
\]
or any other word with a nonterminal boundary symbol as admissible proof-side inputs.
\end{corollary}

\begin{proof}
Immediate from Proposition~\ref{prop:certified-extension-words-have-exactly-form-Ak-or-AkB}.
\end{proof}

\begin{remark}[What this word form now buys]
\label{rem:what-this-word-form-now-buys}
This is the clean point where the active line becomes readable by a genuinely small working alphabet. The theorematic side still decides each local stage, but once those decisions are certified, later layers may consume the finite word \(\mathbf w\) rather than the full local proof tree. That is the part of the later optimization discussion worth keeping: decoder maps, QAM readout, Markov summaries, Monte-Carlo deviation profiles, and matrix-side diagnostics may all work over the certified terminal language \(\mathcal L_{\mathrm{ext}}^{\mathrm{term}}\). What they may not do is inject inadmissible words or continue reading letters after an honest terminal \(\mathsf B\)-symbol has already appeared.
\end{remark}



\noindent\textbf{Reader cue.} Read the next block as the certificate-aware downstream matrix interpretation of the certified extension language. Its purpose is not to let the Heisenberg/Millennium side decide theorematic extension questions on its own, but to force every later matrix-side readout to distinguish sharply between certified horizon-growth and certified earliest-boundary exposure.

\subsection*{19AZ$^{(6ab)}$. Certificate-aware horizon/boundary interpretation for Heisenberg/Millennium matrix readout}
\label{sec:certificate-aware-horizon-boundary-interpretation-for-heisenberg-millennium-matrix-readout}
After 19AZ$^{(6aa)}$, the active line carries a small certified extension language. The natural next use of that language is a downstream matrix-side interpretation: the Heisenberg/Millennium layer should no longer be read as an undifferentiated global diagnostic, but as a certificate-aware readout whose outputs are tagged by whether the active theorematic line has presently certified horizon continuation or exposed the earliest honest boundary stage. This keeps the matrix layer useful without allowing it to outrun the proof line.

\begin{definition}[Certificate-aware Heisenberg/Millennium readout tags and readout map]
\label{def:certificate-aware-heisenberg-millennium-readout-tags-and-readout-map}
Fix a current certified horizon \(\mathbf K_m^{\mathrm{cur}}\), and let \(\mathbf w\in \mathcal L_{\mathrm{ext}}^{\mathrm{term}}\) be a certified finite extension word in the sense of Corollary~\ref{cor:extension-protocol-defines-regular-terminal-word-language}. Introduce the two matrix-side readout tags
\[
  \mathsf{HZ}_{m}(\mathbf w) = \text{``horizon-certified matrix readout on the observed range''},
\]
\[
  \mathsf{BD}_{m}(\mathbf w) = \text{``boundary-certified matrix readout on the observed range''}.
\]
Define the certificate-aware downstream readout map
\[
  \mathsf{Read}^{\mathrm{MM/HC}}_{m}(\mathbf w)
  \in
  \{\mathsf{HZ}_{m}(\mathbf w),\mathsf{BD}_{m}(\mathbf w)\}
\]
by
\[
  \mathsf{Read}^{\mathrm{MM/HC}}_{m}(\mathbf w)
  =
  \begin{cases}
    \mathsf{HZ}_{m}(\mathbf w), & \text{if }\mathbf w=\mathsf A^k\text{ for some }k\ge 0,\\[2mm]
    \mathsf{BD}_{m}(\mathbf w), & \text{if }\mathbf w=\mathsf A^k\mathsf B\text{ for some }k\ge 0.
  \end{cases}
\]
Interpretation:
\begin{itemize}[leftmargin=2em]
  \item \(\mathsf{HZ}_{m}(\mathbf w)\) means that, on the observed certified range, the matrix side is only licensed to read a horizon-persistent transport picture;
  \item \(\mathsf{BD}_{m}(\mathbf w)\) means that, on the observed certified range, the matrix side must switch to earliest-boundary exposure and may not read past the terminal boundary stage as though further certified transport had already been proved.
\end{itemize}
\end{definition}

\begin{proposition}[Certificate-aware matrix readout is faithful to the certified extension language]
\label{prop:certificate-aware-matrix-readout-is-faithful-to-certified-extension-language}
Assume the setting of Definition~\ref{def:certificate-aware-heisenberg-millennium-readout-tags-and-readout-map}. Then:
\begin{enumerate}[label=(\roman*)]
  \item the readout map \(\mathsf{Read}^{\mathrm{MM/HC}}_{m}\) is well-defined on every certified finite extension word;
  \item the matrix-side tag distinguishes exactly the two proof-legitimate cases ``advance-only on the observed range'' and ``advance-until-earliest-boundary on the observed range'';
  \item no downstream Heisenberg/Millennium interpretation is licensed to convert a boundary-certified word \(\mathsf A^k\mathsf B\) into a horizon-certified continuation reading, and no horizon-certified word \(\mathsf A^k\) is allowed to be advertised as exposing an already certified boundary event that the theorematic line has not produced.
\end{enumerate}
\end{proposition}

\begin{proof}
By Corollary~\ref{cor:extension-protocol-defines-regular-terminal-word-language}, every certified finite extension word has one of the two forms \(\mathsf A^k\) or \(\mathsf A^k\mathsf B\). The definition therefore assigns exactly one matrix-side tag to every certified word, so the readout map is well-defined, proving (i). Since these are precisely the two admissible word forms, the map distinguishes exactly the two proof-legitimate cases listed in (ii). Clause (iii) is then just the statement that the downstream matrix layer must remain faithful to the certified word it receives: a terminal \(\mathsf B\)-symbol cannot be reinterpreted as if certified continuation had already occurred, and the absence of a terminal \(\mathsf B\)-symbol cannot be inflated into a claimed certified boundary exposure.
\end{proof}

\begin{corollary}[The present \(\sigma_3\) package already induces a first certificate-aware matrix reading]
\label{cor:present-sigma3-package-already-induces-first-certificate-aware-matrix-reading}
In the current active line, the certified first-stage package from 19AZ$^{(6v)}$--19AZ$^{(6aa)}$ already induces a first certificate-aware Heisenberg/Millennium readout. Concretely, the first-stage word is either
\[
  \mathbf w_3=\mathsf A
  \qquad\text{or}\qquad
  \mathbf w_3=\mathsf B,
\]
and therefore the associated downstream matrix-side reading is already forced into exactly one of the two tags
\[
  \mathsf{Read}^{\mathrm{MM/HC}}_{2}(\mathbf w_3)
  \in
  \{\mathsf{HZ}_{2}(\mathbf w_3),\mathsf{BD}_{2}(\mathbf w_3)\}.
\]
Thus the matrix side can now be asked a precise question: does the current certified stage read as horizon-persistent transport, or as earliest-boundary exposure?
\end{corollary}

\begin{proof}
Immediate from the first-stage encoding in 19AZ$^{(6v)}$--19AZ$^{(6aa)}$ together with Corollary~\ref{cor:extension-protocol-defines-regular-terminal-word-language} and Proposition~\ref{prop:certificate-aware-matrix-readout-is-faithful-to-certified-extension-language}.
\end{proof}

\begin{remark}[What this matrix reading now buys for later crystal and field-equation upgrades]
\label{rem:what-this-matrix-reading-now-buys-for-later-crystal-and-field-equation-upgrades}
This is the clean interface the later upgrade discussion was missing. The Heisenberg/Millennium layer can now be used as a certificate-aware reader of the active extension state rather than as a competing source of theorematic authority. In particular, later crystal-layer refinements and later field-equation diagnostics should be formulated as refinements of either the horizon-certified readout tag or the boundary-certified readout tag. That means they may sharpen what the current certified state looks like, but they should not overrule whether the active line currently sits in a horizon-growth state or at an earliest-boundary exposure state. The same discipline keeps Monte-Carlo or decoder-side diagnostics honest: they may profile, compress, or compare the tagged state, but not replace the proof-side tag assignment itself.
\end{remark}


\noindent\textbf{Reader cue.} Read the next block as the certificate-aware downstream crystal interpretation of the matrix-side tag discipline. Its purpose is not to reopen the earlier addressed-crystal architecture at full generality, but to force every later crystal-side readout to distinguish sharply between horizon-persistent transport/domain structure and earliest-boundary defect exposure.

\subsection*{19AZ$^{(6ac)}$. Certificate-aware horizon/boundary interpretation for crystal-layer readout}
\label{sec:certificate-aware-horizon-boundary-interpretation-for-crystal-layer-readout}
After 19AZ$^{(6ab)}$, the active line no longer carries only a proof-side token language and a matrix-side tag. It also carries the first clean rule for how the later crystal layer is allowed to read the active extension state. The crystal side should therefore not be treated as a free global metaphor or as an independent decision engine. It should be read as a downstream structure-preserving refinement of the already certified horizon-versus-boundary status.

\begin{definition}[Certificate-aware crystal readout tags and crystal-side readout map]
\label{def:certificate-aware-crystal-readout-tags-and-crystal-side-readout-map}
Fix a current certified horizon \(\mathbf K_m^{\mathrm{cur}}\), let \(\mathbf w\in \mathcal L_{\mathrm{ext}}^{\mathrm{term}}\) be a certified finite extension word, and let
\[
  \mathsf{Read}^{\mathrm{MM/HC}}_{m}(\mathbf w)
  \in
  \{\mathsf{HZ}_{m}(\mathbf w),\mathsf{BD}_{m}(\mathbf w)\}
\]
be the associated certificate-aware matrix-side tag from Definition~\ref{def:certificate-aware-heisenberg-millennium-readout-tags-and-readout-map}. Introduce the two crystal-side readout tags
\[
  \mathsf{CRH}_{m}(\mathbf w)=\text{``horizon-certified crystal/domain transport readout''},
\]
\[
  \mathsf{CRB}_{m}(\mathbf w)=\text{``boundary-certified crystal defect/front readout''}.
\]
Define the certificate-aware crystal-side readout map
\[
  \mathsf{Read}^{\mathrm{cr}}_{m}(\mathbf w)
  \in
  \{\mathsf{CRH}_{m}(\mathbf w),\mathsf{CRB}_{m}(\mathbf w)\}
\]
by
\[
  \mathsf{Read}^{\mathrm{cr}}_{m}(\mathbf w)
  =
  \begin{cases}
    \mathsf{CRH}_{m}(\mathbf w), & \text{if }\mathsf{Read}^{\mathrm{MM/HC}}_{m}(\mathbf w)=\mathsf{HZ}_{m}(\mathbf w),\\[2mm]
    \mathsf{CRB}_{m}(\mathbf w), & \text{if }\mathsf{Read}^{\mathrm{MM/HC}}_{m}(\mathbf w)=\mathsf{BD}_{m}(\mathbf w).
  \end{cases}
\]
Interpretation:
\begin{itemize}[leftmargin=2em]
  \item \(\mathsf{CRH}_{m}(\mathbf w)\) means that the crystal side is licensed to read the currently certified range as one transport-persistent domain picture, with no theorematically certified defect-front beyond that range yet exposed;
  \item \(\mathsf{CRB}_{m}(\mathbf w)\) means that the crystal side must treat the terminal certified information as an earliest honest defect/front exposure and may not smooth, continue, or visually homogenize that boundary as though later corridor coherence had already been proved.
\end{itemize}
\end{definition}

\begin{proposition}[Certificate-aware crystal readout is faithful to the matrix-side horizon/boundary tag]
\label{prop:certificate-aware-crystal-readout-is-faithful-to-the-matrix-side-horizon-boundary-tag}
Assume the setting of Definition~\ref{def:certificate-aware-crystal-readout-tags-and-crystal-side-readout-map}. Then:
\begin{enumerate}[label=(\roman*)]
  \item the crystal-side readout map \(\mathsf{Read}^{\mathrm{cr}}_{m}\) is well-defined on every certified finite extension word;
  \item the crystal-side tag distinguishes exactly the two proof-legitimate crystal readings ``domain transport on the certified range'' and ``earliest defect/front exposure at the certified boundary stage'';
  \item no downstream crystal-layer interpretation is licensed to smooth a boundary-certified state into a transport-certified domain picture, and no horizon-certified state is allowed to be advertised as already carrying a proved defect/front event that the active line has not certified.
\end{enumerate}
\end{proposition}

\begin{proof}
By Proposition~\ref{prop:certificate-aware-matrix-readout-is-faithful-to-certified-extension-language}, the matrix-side readout assigns exactly one of the two tags \(\mathsf{HZ}_{m}(\mathbf w)\) or \(\mathsf{BD}_{m}(\mathbf w)\) to every certified word. Definition~\ref{def:certificate-aware-crystal-readout-tags-and-crystal-side-readout-map} therefore assigns exactly one crystal-side tag \(\mathsf{CRH}_{m}(\mathbf w)\) or \(\mathsf{CRB}_{m}(\mathbf w)\), proving (i). Since these two tags simply refine the two proof-legitimate matrix-side statuses into crystal/domain language, statement (ii) follows immediately. Clause (iii) is the same faithfulness discipline expressed on the crystal side: a boundary-certified terminal stage cannot be visually or narratively turned into uninterrupted continuation, and a transport-certified range cannot be embellished with a certified defect/front event that the theorematic line has not produced.
\end{proof}

\begin{corollary}[The present \(\sigma_3\) package already induces a first certificate-aware crystal reading]
\label{cor:present-sigma3-package-already-induces-first-certificate-aware-crystal-reading}
In the current active line, the first-stage word \(\mathbf w_3\in\{\mathsf A,\mathsf B\}\) already induces a first certificate-aware crystal-side readout
\[
  \mathsf{Read}^{\mathrm{cr}}_{2}(\mathbf w_3)
  \in
  \{\mathsf{CRH}_{2}(\mathbf w_3),\mathsf{CRB}_{2}(\mathbf w_3)\}.
\]
Thus the crystal layer can now be asked a precise and conservative question: does the present certified stage still read as one transport-persistent domain picture, or has the earliest honest defect/front already been exposed at the current terminal stage?
\end{corollary}

\begin{proof}
Immediate from Corollary~\ref{cor:present-sigma3-package-already-induces-first-certificate-aware-matrix-reading} and Proposition~\ref{prop:certificate-aware-crystal-readout-is-faithful-to-the-matrix-side-horizon-boundary-tag}.
\end{proof}

\begin{remark}[What this crystal reading now buys]
\label{rem:what-this-crystal-reading-now-buys}
This is the clean point where the earlier liquid-crystal/domain/defect language becomes structurally disciplined by the active extension certificates. The crystal side may now refine the current state into domain persistence, defect loading, or front exposure language without pretending to certify a later slot on its own. In particular, later holographic boundary visuals, domain-defect compression schemes, and Monte-Carlo or decoder-side pattern scans may profile the crystal-side tag, but they must inherit it from the certified extension language rather than invent it.
\end{remark}

\noindent\textbf{Reader cue.} Read the next block as the certificate-aware downstream pre-equation interface for field-equation diagnostics. Its purpose is not to claim a finished Einstein-, Dirac-, or QFT-type projection law at this stage, but to force every later field-equation ansatz to distinguish sharply between a horizon-certified continuation profile and an earliest-boundary source/exposure profile.

\subsection*{19AZ$^{(6ad)}$. Certificate-aware horizon/boundary interpretation for field-equation diagnostics}
\label{sec:certificate-aware-horizon-boundary-interpretation-for-field-equation-diagnostics}
After 19AZ$^{(6ac)}$, the active line has a theorematically certified extension language, a matrix-side readout, and a crystal-side readout. The next honest compression step is not yet a full new field equation. It is a certificate-aware pre-equation interface that tells any later physical or PDE-style readout whether it is looking at a continuation-certified transport profile or at an earliest-boundary exposure profile. This is the smallest disciplined upgrade that the current line can honestly support.

\begin{definition}[Certificate-aware field-diagnostic tags and pre-equation map]
\label{def:certificate-aware-field-diagnostic-tags-and-pre-equation-map}
Fix a current certified horizon \(\mathbf K_m^{\mathrm{cur}}\), let \(\mathbf w\in \mathcal L_{\mathrm{ext}}^{\mathrm{term}}\) be a certified finite extension word, and let
\[
  \mathsf{Read}^{\mathrm{cr}}_{m}(\mathbf w)
  \in
  \{\mathsf{CRH}_{m}(\mathbf w),\mathsf{CRB}_{m}(\mathbf w)\}
\]
be the associated crystal-side readout from Definition~\ref{def:certificate-aware-crystal-readout-tags-and-crystal-side-readout-map}. Introduce the two field-diagnostic tags
\[
  \mathsf{FEH}_{m}(\mathbf w)=\text{``horizon-certified pre-equation transport profile''},
\]
\[
  \mathsf{FEB}_{m}(\mathbf w)=\text{``boundary-certified pre-equation source/exposure profile''}.
\]
Define the certificate-aware field-diagnostic map
\[
  \mathsf{Diag}^{\mathrm{field}}_{m}(\mathbf w)
  \in
  \{\mathsf{FEH}_{m}(\mathbf w),\mathsf{FEB}_{m}(\mathbf w)\}
\]
by
\[
  \mathsf{Diag}^{\mathrm{field}}_{m}(\mathbf w)
  =
  \begin{cases}
    \mathsf{FEH}_{m}(\mathbf w), & \text{if }\mathsf{Read}^{\mathrm{cr}}_{m}(\mathbf w)=\mathsf{CRH}_{m}(\mathbf w),\\[2mm]
    \mathsf{FEB}_{m}(\mathbf w), & \text{if }\mathsf{Read}^{\mathrm{cr}}_{m}(\mathbf w)=\mathsf{CRB}_{m}(\mathbf w).
  \end{cases}
\]
Interpretation:
\begin{itemize}[leftmargin=2em]
  \item \(\mathsf{FEH}_{m}(\mathbf w)\) means that any later field-equation ansatz is presently licensed only to read continuation-compatible transport on the already certified range;
  \item \(\mathsf{FEB}_{m}(\mathbf w)\) means that any later field-equation ansatz must already acknowledge an earliest honest boundary/source exposure at the terminal certified stage and may not write a smooth forward continuation law past that stage as though the proof line had already certified it.
\end{itemize}
\end{definition}

\begin{proposition}[Certificate-aware field diagnostics are faithful to the crystal-side tag]
\label{prop:certificate-aware-field-diagnostics-are-faithful-to-the-crystal-side-tag}
Assume the setting of Definition~\ref{def:certificate-aware-field-diagnostic-tags-and-pre-equation-map}. Then:
\begin{enumerate}[label=(\roman*)]
  \item the field-diagnostic map \(\mathsf{Diag}^{\mathrm{field}}_{m}\) is well-defined on every certified finite extension word;
  \item the field-diagnostic tag distinguishes exactly the two presently legitimate pre-equation situations ``continuation-compatible certified transport'' and ``earliest-boundary source/exposure already active'';
  \item no downstream field-equation ansatz is licensed to erase a terminal boundary-certified status, and no continuation-certified status is allowed to be sold as a proved boundary/source event that the active line has not certified.
\end{enumerate}
\end{proposition}

\begin{proof}
By Proposition~\ref{prop:certificate-aware-crystal-readout-is-faithful-to-the-matrix-side-horizon-boundary-tag}, the crystal-side readout assigns exactly one of the tags \(\mathsf{CRH}_{m}(\mathbf w)\) or \(\mathsf{CRB}_{m}(\mathbf w)\) to every certified word. Definition~\ref{def:certificate-aware-field-diagnostic-tags-and-pre-equation-map} therefore assigns exactly one field-diagnostic tag \(\mathsf{FEH}_{m}(\mathbf w)\) or \(\mathsf{FEB}_{m}(\mathbf w)\), proving (i). Since these are precisely the two downstream field-diagnostic refinements of the crystal-side statuses, statement (ii) follows. Clause (iii) is the same certificate faithfulness expressed at the pre-equation layer: a certified terminal boundary/source exposure cannot be rewritten as already-proved smooth continuation, and a certified transport range cannot be exaggerated into a proved boundary event.
\end{proof}

\begin{corollary}[The present \(\sigma_3\) package already induces a first certificate-aware field-diagnostic split]
\label{cor:present-sigma3-package-already-induces-first-certificate-aware-field-diagnostic-split}
In the current active line, the first-stage word \(\mathbf w_3\in\{\mathsf A,\mathsf B\}\) already induces a first certificate-aware field-diagnostic split
\[
  \mathsf{Diag}^{\mathrm{field}}_{2}(\mathbf w_3)
  \in
  \{\mathsf{FEH}_{2}(\mathbf w_3),\mathsf{FEB}_{2}(\mathbf w_3)\}.
\]
Thus the later field-equation layer can already be asked a disciplined question: is the current certified stage to be read as continuation-compatible transport on the observed range, or as an earliest-boundary source/exposure profile that must remain explicit in every later PDE or operator ansatz?
\end{corollary}

\begin{proof}
Immediate from Corollary~\ref{cor:present-sigma3-package-already-induces-first-certificate-aware-crystal-reading} and Proposition~\ref{prop:certificate-aware-field-diagnostics-are-faithful-to-the-crystal-side-tag}.
\end{proof}

\begin{remark}[What this field-diagnostic interface now buys]
\label{rem:what-this-field-diagnostic-interface-now-buys}
This is the conservative upgrade the later field-equation discussion was missing. The active line still does not claim a finished physical PDE, a finished Einstein equation, or a finished emergence law. What it now does claim is a disciplined pre-equation interface: every later ansatz must declare whether it is reading a horizon-certified transport state or an earliest-boundary source/exposure state. That already sharpens admissibility. A later field equation is acceptable only if its local readout remains faithful to the certified tag it inherits from the extension language.
\end{remark}


\noindent\textbf{Reader cue.} Read the next block as the admissibility and consistency gate for any later PDE-, operator-, or emergence-style ansatz built on the certificate-aware field-diagnostic interface. Its purpose is not to force one unique final field equation, but to exclude exactly those later ansätze that would erase, overstate, or misroute the certified horizon/boundary status already inherited from the active extension line.

\subsection*{19AZ$^{(6ae)}$. Admissibility and consistency conditions for later field-equation ansätze}
\label{sec:admissibility-and-consistency-conditions-for-later-field-equation-ansatze}
After 19AZ$^{(6ad)}$, the active line still does not claim a unique finished field equation. What it does claim is stronger than a vague heuristic: any later PDE-, operator-, or emergence-style law must now pass a certificate-aware admissibility gate. In other words, once a later ansatz inherits the field-diagnostic tag \(\mathsf{FEH}_{m}(\mathbf w)\) or \(\mathsf{FEB}_{m}(\mathbf w)\), it is no longer free to ignore the certified status that generated that tag.

\begin{definition}[Certificate-aware admissibility for later field-equation ansätze]
\label{def:certificate-aware-admissibility-for-later-field-equation-ansatze}
Fix a current certified horizon \(\mathbf K_m^{\mathrm{cur}}\), let \(\mathbf w\in \mathcal L_{\mathrm{ext}}^{\mathrm{term}}\) be a certified finite extension word, and let
\[
  \mathsf{Diag}^{\mathrm{field}}_{m}(\mathbf w)
  \in
  \{\mathsf{FEH}_{m}(\mathbf w),\mathsf{FEB}_{m}(\mathbf w)\}
\]
be the inherited field-diagnostic tag from Definition~\ref{def:certificate-aware-field-diagnostic-tags-and-pre-equation-map}. A later PDE-, operator-, transport-, or emergence-style ansatz
\[
  \mathfrak E_{m,\mathbf w}
\]
is called \emph{certificate-aware admissible} if it satisfies the following conditions.
\begin{enumerate}[label=(\roman*)]
  \item \textbf{Tag fidelity.} The ansatz must explicitly declare whether it is reading the state under \(\mathsf{FEH}_{m}(\mathbf w)\) or under \(\mathsf{FEB}_{m}(\mathbf w)\), and its local interpretation must respect that inherited tag.
  \item \textbf{No boundary erasure.} If the inherited tag is \(\mathsf{FEB}_{m}(\mathbf w)\), the ansatz may not replace the terminal certified boundary/source exposure by a smooth forward continuation law as though the active line had already certified later corridor transport.
  \item \textbf{No boundary inflation.} If the inherited tag is \(\mathsf{FEH}_{m}(\mathbf w)\), the ansatz may not advertise an already certified source, defect, singularity, or terminal boundary event that the active line has not produced.
  \item \textbf{Range discipline.} The ansatz may refine, profile, or parameterize the already certified observed range, but it may not claim theorematic authority beyond the range encoded by the certified word \(\mathbf w\).
\end{enumerate}
If one of these conditions fails, the ansatz is called \emph{certificate-aware inadmissible} on the input pair \((m,\mathbf w)\).
\end{definition}

\begin{proposition}[The certificate-aware field interface induces a sharp admissibility dichotomy]
\label{prop:certificate-aware-field-interface-induces-sharp-admissibility-dichotomy}
Assume the setting of Definition~\ref{def:certificate-aware-admissibility-for-later-field-equation-ansatze}. Then every later field-equation ansatz \(\mathfrak E_{m,\mathbf w}\) is forced into exactly one of the two statuses:
\[
  \text{certificate-aware admissible}
  \qquad\text{or}\qquad
  \text{certificate-aware inadmissible}.
\]
Moreover:
\begin{enumerate}[label=(\roman*)]
  \item every admissible ansatz is a downstream refinement of the inherited tag rather than a replacement for it;
  \item every inadmissible ansatz fails by one of exactly three routes: boundary erasure, boundary inflation, or range overreach;
  \item this admissibility test is independent of whether the later ansatz is written in matrix, crystal, PDE, operator, Monte-Carlo-assisted, or decoder-compressed language.
\end{enumerate}
\end{proposition}

\begin{proof}
Definition~\ref{def:certificate-aware-admissibility-for-later-field-equation-ansatze} gives explicit necessary and sufficient conditions for admissibility on the input pair \((m,\mathbf w)\). Hence every later ansatz either satisfies all of them or fails at least one, which proves the displayed dichotomy. Statement (i) is immediate from tag fidelity together with range discipline: an admissible ansatz inherits the certified field tag and may only refine the already certified state. For (ii), the only ways to violate the conditions are exactly the ones listed in the definition: either an inherited boundary-certified state is erased, an inherited continuation-certified state is exaggerated into a certified boundary/source event, or the ansatz claims interpretive authority beyond the certified observed range. Statement (iii) follows because the admissibility test depends only on the inherited certified tag and on faithfulness to the certified word, not on the notation or computational packaging chosen downstream.
\end{proof}

\begin{corollary}[The present \(\sigma_3\) package already excludes the wrong kinds of field-law upgrades]
\label{cor:present-sigma3-package-already-excludes-wrong-kinds-of-field-law-upgrades}
In the current active line, the first-stage package with \(\mathbf w_3\in\{\mathsf A,\mathsf B\}\) already excludes two opposite classes of later upgrade errors:
\begin{enumerate}[label=(\roman*)]
  \item if \(\mathbf w_3=\mathsf B\), no later matrix-, crystal-, or field-law ansatz is allowed to smooth that first certified boundary exposure into an already proved forward continuation picture;
  \item if \(\mathbf w_3=\mathsf A\), no later matrix-, crystal-, or field-law ansatz is allowed to present the current stage as though an already certified boundary/source event had been exposed.
\end{enumerate}
Thus even before a finished field equation exists, the present line already knows which later upgrade moves are structurally forbidden.
\end{corollary}

\begin{proof}
Immediate from Corollary~\ref{cor:present-sigma3-package-already-induces-first-certificate-aware-field-diagnostic-split} and Proposition~\ref{prop:certificate-aware-field-interface-induces-sharp-admissibility-dichotomy}.
\end{proof}

\begin{remark}[What this admissibility gate now buys]
\label{rem:what-this-admissibility-gate-now-buys}
This is the point where the current line becomes strong enough to police its own future upgrades. A later Einstein-style, transport-style, source-law, Monte-Carlo-assisted closure law, or operator-valued emergence ansatz no longer has to be rejected because it is ambitious; it has to be rejected only if it violates the inherited certificate. That is exactly the right level of rigidity at the present stage: permissive about representation, strict about structural faithfulness.
\end{remark}

\noindent\textbf{Reader cue.} Read the next block as the certificate-aware upgrade-routing interface that tells every later refinement layer which type of work is legitimately allowed after a field-diagnostic tag has been inherited. Its purpose is not to force one unique downstream model, but to route conservative continuation-style upgrades to the horizon-certified branch and boundary-profiling upgrades to the boundary-certified branch, while excluding cross-routings that would reopen or falsify the certified status.

\subsection*{19AZ$^{(6af)}$. Upgrade-routing consequences of certificate-aware field-diagnostic tags}
\label{sec:upgrade-routing-consequences-of-certificate-aware-field-diagnostic-tags}
After 19AZ$^{(6ae)}$, the active line has more than a static admissibility gate. It now has enough structure to route later upgrades according to the inherited certificate. The point is simple: not every technically sophisticated downstream refinement is appropriate for every certified word. A horizon-certified tag should route later work toward conservative continuation-style refinement on the already proved range, while a boundary-certified tag should route later work toward earliest-boundary profiling, exposure analysis, and stop-sensitive diagnostics. This routing discipline is the missing interface between the theorematic extension line and later matrix-, crystal-, field-, decoder-, or probe-based elaborations.

\begin{definition}[Certificate-aware upgrade-routing map]
\label{def:certificate-aware-upgrade-routing-map}
Fix a current certified horizon \(\mathbf K_m^{\mathrm{cur}}\), let \(\mathbf w\in \mathcal L_{\mathrm{ext}}^{\mathrm{term}}\) be a certified finite extension word, and let
\[
  \mathsf{Diag}^{\mathrm{field}}_{m}(\mathbf w)
  \in
  \{\mathsf{FEH}_{m}(\mathbf w),\mathsf{FEB}_{m}(\mathbf w)\}
\]
be the inherited field-diagnostic tag from Definition~\ref{def:certificate-aware-field-diagnostic-tags-and-pre-equation-map}. Introduce the two conservative upgrade-route classes
\begin{align*}
  \mathcal U^{\mathrm{hor}}_{m}(\mathbf w)
  &:= \text{``continuation-compatible refinement on the already certified range''},\\
  \mathcal U^{\mathrm{bdry}}_{m}(\mathbf w)
  &:= \text{``earliest-boundary profiling and stop-sensitive exposure analysis''}.
\end{align*}
Define the certificate-aware upgrade-routing map
\[
  \mathsf{Route}^{\mathrm{upg}}_{m}(\mathbf w)
  \in
  \{\mathcal U^{\mathrm{hor}}_{m}(\mathbf w),\mathcal U^{\mathrm{bdry}}_{m}(\mathbf w)\}
\]
by
\[
  \mathsf{Route}^{\mathrm{upg}}_{m}(\mathbf w)
  =
  \begin{cases}
    \mathcal U^{\mathrm{hor}}_{m}(\mathbf w), & \text{if }\mathsf{Diag}^{\mathrm{field}}_{m}(\mathbf w)=\mathsf{FEH}_{m}(\mathbf w),\\[2mm]
    \mathcal U^{\mathrm{bdry}}_{m}(\mathbf w), & \text{if }\mathsf{Diag}^{\mathrm{field}}_{m}(\mathbf w)=\mathsf{FEB}_{m}(\mathbf w).
  \end{cases}
\]
Interpretation:
\begin{itemize}[leftmargin=2em]
  \item the horizon-certified route \(\mathcal U^{\mathrm{hor}}_{m}(\mathbf w)\) licenses only those later upgrades that refine transport, matrix readout, crystal stabilization, or PDE/operator behavior on the range already certified by \(\mathbf w\), without claiming a new boundary/source event;
  \item the boundary-certified route \(\mathcal U^{\mathrm{bdry}}_{m}(\mathbf w)\) licenses only those later upgrades that sharpen the earliest-boundary profile itself, for example by defect localization, source/exposure analysis, or stop-sensitive decoder/probe diagnostics, without pretending that corridor continuation past the terminal boundary has already been proved;
  \item any downstream refinement that crosses these route classes against the inherited tag is called a \emph{cross-routed upgrade} and is provisionally inadmissible until a new local theorematic stage changes the certified status.
\end{itemize}
\end{definition}

\begin{proposition}[The inherited field-diagnostic tag forces a unique conservative upgrade route]
\label{prop:inherited-field-diagnostic-tag-forces-unique-conservative-upgrade-route}
Assume the setting of Definition~\ref{def:certificate-aware-upgrade-routing-map}. Then:
\begin{enumerate}[label=(\roman*)]
  \item the routing map \(\mathsf{Route}^{\mathrm{upg}}_{m}\) is well-defined on every certified finite extension word;
  \item a horizon-certified tag \(\mathsf{FEH}_{m}(\mathbf w)\) forces the conservative route class \(\mathcal U^{\mathrm{hor}}_{m}(\mathbf w)\), while a boundary-certified tag \(\mathsf{FEB}_{m}(\mathbf w)\) forces the conservative route class \(\mathcal U^{\mathrm{bdry}}_{m}(\mathbf w)\);
  \item every cross-routed upgrade is incompatible with the certificate-aware admissibility conditions from Definition~\ref{def:certificate-aware-admissibility-for-later-field-equation-ansatze} unless and until a new theorematic local stage certifies a different status.
\end{enumerate}
\end{proposition}

\begin{proof}
By Proposition~\ref{prop:certificate-aware-field-interface-induces-sharp-admissibility-dichotomy}, every certified word inherits exactly one field-diagnostic tag, either \(\mathsf{FEH}_{m}(\mathbf w)\) or \(\mathsf{FEB}_{m}(\mathbf w)\). Definition~\ref{def:certificate-aware-upgrade-routing-map} therefore assigns exactly one conservative route class, proving (i). Clause (ii) is the explicit content of the definition. For (iii), suppose a downstream refinement is routed against the inherited tag. If a boundary-certified word is forced onto the continuation-style route, then the upgrade attempts boundary erasure by treating an earliest-boundary exposure as already-certified transport continuation. If a horizon-certified word is forced onto the boundary-style route, then the upgrade attempts boundary inflation by advertising a defect, source, or stopping event not certified by the active line. Both behaviors violate the admissibility gate from Definition~\ref{def:certificate-aware-admissibility-for-later-field-equation-ansatze}. Hence cross-routed upgrades are incompatible with the present certificate unless a later local theorematic stage genuinely changes that certificate.
\end{proof}

\begin{corollary}[The present \(\sigma_3\) package already induces a first conservative upgrade-routing split]
\label{cor:present-sigma3-package-already-induces-first-conservative-upgrade-routing-split}
In the current active line, the first-stage word \(\mathbf w_3\in\{\mathsf A,\mathsf B\}\) already induces a first conservative routing split
\[
  \mathsf{Route}^{\mathrm{upg}}_{2}(\mathbf w_3)
  \in
  \{\mathcal U^{\mathrm{hor}}_{2}(\mathbf w_3),\mathcal U^{\mathrm{bdry}}_{2}(\mathbf w_3)\}.
\]
Thus the next downstream question is no longer vague. If the inherited tag is horizon-certified, later work should concentrate on transport-compatible refinement of the certified range. If the inherited tag is boundary-certified, later work should concentrate on profiling the earliest honest boundary rather than pretending that the corridor has already been extended through it.
\end{corollary}

\begin{proof}
Immediate from Corollary~\ref{cor:present-sigma3-package-already-induces-first-certificate-aware-field-diagnostic-split} and Proposition~\ref{prop:inherited-field-diagnostic-tag-forces-unique-conservative-upgrade-route}.
\end{proof}

\begin{remark}[What this upgrade-routing interface now buys]
\label{rem:what-this-upgrade-routing-interface-now-buys}
This is the first place where the active line tells later heuristic or computational layers not only what they may \emph{say}, but also what they should \emph{do next}. Matrix sharpening, crystal stabilization, conservative PDE refinement, decoder compression, Monte-Carlo probing, or source/exposure diagnostics are now routed by the inherited certificate rather than by taste. In particular, Monte-Carlo, decoder, hash, or Markov-style probes may read and compress the routed status, but they are not allowed to replace the theorematic routing decision itself. That keeps the later tool layer useful without allowing it to overrule the certified extension line.
\end{remark}




\noindent\textbf{Reader cue.} Read the next block as the normalized follow-on upgrade menu sitting \emph{below} the certificate-aware routing interface. Its purpose is to say which concrete families of downstream refinement are genuinely licensed once a route has already been fixed, and which families remain out of bounds until a later local theorematic stage changes the certificate.

\subsection*{19AZ$^{(6ag)}$. Normalized follow-on upgrade families below the certificate-aware routing interface}
\label{sec:normalized-follow-on-upgrade-families-below-the-certificate-aware-routing-interface}
After 19AZ$^{(6af)}$, the line knows not only whether a later refinement is admissible, but also which route it must follow. The next honest compression step is to normalize the concrete follow-on upgrade families that sit below those routes. This is important because later readers should not have to guess whether, for example, matrix sharpening, crystal stabilization, source localization, decoder compression, or Monte-Carlo probing belongs to the continuation side or to the earliest-boundary side. The present block turns that tacit distinction into a finite, certificate-aware menu.

\begin{definition}[Normalized follow-on upgrade families]
\label{def:normalized-follow-on-upgrade-families}
Fix a current certified horizon \(\mathbf K_m^{\mathrm{cur}}\), let \(\mathbf w\in \mathcal L_{\mathrm{ext}}^{\mathrm{term}}\) be a certified finite extension word, and let
\[
  \mathsf{Route}^{\mathrm{upg}}_{m}(\mathbf w)
  \in
  \{\mathcal U^{\mathrm{hor}}_{m}(\mathbf w),\mathcal U^{\mathrm{bdry}}_{m}(\mathbf w)\}
\]
be the inherited upgrade route from Definition~\ref{def:certificate-aware-upgrade-routing-map}.

Define the \emph{horizon-follow-on family menu}
\[
  \mathcal F^{\mathrm{hor}}_{m}(\mathbf w)
  :=
  \{\mathsf{MSH}_{m}(\mathbf w),\mathsf{CRS}_{m}(\mathbf w),\mathsf{PDE}^{\mathrm{cons}}_{m}(\mathbf w),\mathsf{DEC}^{\mathrm{hor}}_{m}(\mathbf w)\}
\]
with the intended meanings
\begin{align*}
  \mathsf{MSH}_{m}(\mathbf w) &:= \text{``matrix sharpening on the already certified range''},\\
  \mathsf{CRS}_{m}(\mathbf w) &:= \text{``crystal stabilization / transport-consistency refinement''},\\
  \mathsf{PDE}^{\mathrm{cons}}_{m}(\mathbf w) &:= \text{``conservative PDE/operator refinement on the certified range''},\\
  \mathsf{DEC}^{\mathrm{hor}}_{m}(\mathbf w) &:= \text{``decoder / compression readout of a horizon-certified status''}.
\end{align*}

Define the \emph{boundary-follow-on family menu}
\[
  \mathcal F^{\mathrm{bdry}}_{m}(\mathbf w)
  :=
  \{\mathsf{SRC}_{m}(\mathbf w),\mathsf{DFX}_{m}(\mathbf w),\mathsf{MC}^{\mathrm{bdry}}_{m}(\mathbf w),\mathsf{DEC}^{\mathrm{bdry}}_{m}(\mathbf w)\}
\]
with the intended meanings
\begin{align*}
  \mathsf{SRC}_{m}(\mathbf w) &:= \text{``source / exposure localization at the earliest honest boundary''},\\
  \mathsf{DFX}_{m}(\mathbf w) &:= \text{``defect-front or crystal-boundary profiling''},\\
  \mathsf{MC}^{\mathrm{bdry}}_{m}(\mathbf w) &:= \text{``stop-sensitive Monte-Carlo or probe diagnostics at the boundary side''},\\
  \mathsf{DEC}^{\mathrm{bdry}}_{m}(\mathbf w) &:= \text{``decoder / compression readout of a boundary-certified status''}.
\end{align*}

A \emph{normalized follow-on upgrade} for \(\mathbf w\) is then any element of the route-compatible menu
\[
  \mathcal F^{\mathrm{follow}}_{m}(\mathbf w)
  :=
  \begin{cases}
    \mathcal F^{\mathrm{hor}}_{m}(\mathbf w), & \text{if }\mathsf{Route}^{\mathrm{upg}}_{m}(\mathbf w)=\mathcal U^{\mathrm{hor}}_{m}(\mathbf w),\\[2mm]
    \mathcal F^{\mathrm{bdry}}_{m}(\mathbf w), & \text{if }\mathsf{Route}^{\mathrm{upg}}_{m}(\mathbf w)=\mathcal U^{\mathrm{bdry}}_{m}(\mathbf w).
  \end{cases}
\]
Any proposed later refinement that selects an element outside \(\mathcal F^{\mathrm{follow}}_{m}(\mathbf w)\) is called a \emph{route-incompatible follow-on upgrade}.
\end{definition}

\begin{proposition}[Every certified word carries a finite route-compatible follow-on upgrade menu]
\label{prop:every-certified-word-carries-a-finite-route-compatible-follow-on-upgrade-menu}
Assume the setting of Definition~\ref{def:normalized-follow-on-upgrade-families}. Then:
\begin{enumerate}[label=(\roman*)]
  \item every certified finite extension word \(\mathbf w\) carries exactly one route-compatible follow-on family menu \(\mathcal F^{\mathrm{follow}}_{m}(\mathbf w)\);
  \item every normalized follow-on upgrade in \(\mathcal F^{\mathrm{follow}}_{m}(\mathbf w)\) preserves the certificate-aware routing discipline from 19AZ$^{(6af)}$;
  \item every route-incompatible follow-on upgrade attempts either continuation overreach on the boundary side or spurious defect/source inflation on the horizon side, and is therefore provisionally inadmissible.
\end{enumerate}
\end{proposition}

\begin{proof}
By Proposition~\ref{prop:inherited-field-diagnostic-tag-forces-unique-conservative-upgrade-route}, every certified finite extension word inherits exactly one conservative upgrade route. Definition~\ref{def:normalized-follow-on-upgrade-families} assigns to that unique route exactly one finite menu, proving (i). Clause (ii) is built into the meanings of the menu entries: the horizon-family items refine only the already certified continuation range, while the boundary-family items profile only the earliest certified boundary or exposure picture. For (iii), if a boundary-certified word is sent to a horizon-family item, the proposed refinement necessarily behaves as though continuation beyond the terminal boundary were already certified; this is continuation overreach. If a horizon-certified word is sent to a boundary-family item, the proposed refinement injects a defect, source, or stopping narrative that has not been certified by the active line; this is spurious inflation. Both violate the route discipline and so are provisionally inadmissible.
\end{proof}

\begin{corollary}[The present \(\sigma_3\) package already carries a finite next-work menu]
\label{cor:present-sigma3-package-already-carries-a-finite-next-work-menu}
In the current active line, the first-stage word \(\mathbf w_3\in\{\mathsf A,\mathsf B\}\) already carries a finite next-work menu
\[
  \mathcal F^{\mathrm{follow}}_{2}(\mathbf w_3)
  \subseteq
  \mathcal F^{\mathrm{hor}}_{2}(\mathbf w_3)\cup \mathcal F^{\mathrm{bdry}}_{2}(\mathbf w_3).
\]
Thus the immediate next work is no longer only route-correct in the abstract. It is already normalized into a finite family of concrete continuation-side or boundary-side upgrade tasks.
\end{corollary}

\begin{proof}
Immediate from Corollary~\ref{cor:present-sigma3-package-already-induces-first-conservative-upgrade-routing-split} and Proposition~\ref{prop:every-certified-word-carries-a-finite-route-compatible-follow-on-upgrade-menu}.
\end{proof}

\begin{remark}[What this finite next-work menu now buys]
\label{rem:what-this-finite-next-work-menu-now-buys}
This is the point where the active line starts to look like a genuine operator's manual for the next minor release. Matrix sharpening, crystal stabilization, conservative PDE refinement, decoder compression, source localization, defect profiling, and stop-sensitive Monte-Carlo probing are no longer a loose wishlist. They have become a finite, certificate-aware menu. In particular, Monte-Carlo now has a safe place: it can act as a boundary-side structural probe when the inherited status is boundary-certified, but it is not allowed to masquerade as theorematic continuation proof.
\end{remark}




oindent	extbf{Reader cue.} Read the next block as the explicit QAM-/decoder-side shadow of the certificate-aware extension interface. Its purpose is not to introduce a new ontology above the current tool layer, but to compress the already inherited field tag, route, and next-work menu into one finite downstream object that QAM- and decoder-style readers may consume without reopening the theorematic proof split.

\subsection*{19AZ$^{(6ah)}$. QAM-/decoder-side shadow of the certificate-aware extension interface}
\label{sec:qam-decoder-side-shadow-of-the-certificate-aware-extension-interface}
After 19AZ$^{(6ag)}$, the active line already carries everything a later QAM- or decoder-side layer needs in order to act conservatively: a certified finite word, a field-diagnostic tag, a forced route, and a finite next-work menu. The next honest compression step is therefore not a new theorem, but a finite shadow object that packages these inherited data into one downstream-readable interface. This is exactly the point where the present certificate-aware tool layer becomes explicitly QAM-/decoder-consumable.

\begin{definition}[QAM-/decoder-side shadow state]
\label{def:qam-decoder-side-shadow-state}
Fix a current certified horizon \(\mathbf K_m^{\mathrm{cur}}\), let \(\mathbf w\in \mathcal L_{\mathrm{ext}}^{\mathrm{term}}\) be a certified finite extension word, and let
\[
  \mathsf{Diag}^{\mathrm{field}}_{m}(\mathbf w)
  \in
  \{\mathsf{FEH}_{m}(\mathbf w),\mathsf{FEB}_{m}(\mathbf w)\},
  \qquad
  \mathsf{Route}^{\mathrm{upg}}_{m}(\mathbf w)
  \in
  \{\mathcal U^{\mathrm{hor}}_{m}(\mathbf w),\mathcal U^{\mathrm{bdry}}_{m}(\mathbf w)\},
\]
and
\[
  \mathcal F^{\mathrm{follow}}_{m}(\mathbf w)
  \in
  \{\mathcal F^{\mathrm{hor}}_{m}(\mathbf w),\mathcal F^{\mathrm{bdry}}_{m}(\mathbf w)\}
\]
be the inherited field tag, upgrade route, and follow-on menu from 19AZ$^{(6ad)}$--19AZ$^{(6ag)}$.

Define the two finite QAM-shadow states
\begin{align*}
  \mathsf{QSH}^{\mathrm{hor}}_{m}(\mathbf w)
  &:=
  \bigl\langle
    \mathsf{FEH}_{m}(\mathbf w),
    \mathcal U^{\mathrm{hor}}_{m}(\mathbf w),
    \mathcal F^{\mathrm{hor}}_{m}(\mathbf w)
  \bigr\rangle,\\
  \mathsf{QSH}^{\mathrm{bdry}}_{m}(\mathbf w)
  &:=
  \bigl\langle
    \mathsf{FEB}_{m}(\mathbf w),
    \mathcal U^{\mathrm{bdry}}_{m}(\mathbf w),
    \mathcal F^{\mathrm{bdry}}_{m}(\mathbf w)
  \bigr\rangle.
\end{align*}
Define the \emph{QAM-/decoder-side shadow map}
\[
  \mathsf{Shadow}^{\mathrm{QAM}}_{m}(\mathbf w)
  \in
  \{\mathsf{QSH}^{\mathrm{hor}}_{m}(\mathbf w),\mathsf{QSH}^{\mathrm{bdry}}_{m}(\mathbf w)\}
\]
by
\[
  \mathsf{Shadow}^{\mathrm{QAM}}_{m}(\mathbf w)
  =
  \begin{cases}
    \mathsf{QSH}^{\mathrm{hor}}_{m}(\mathbf w), & \text{if }\mathsf{Diag}^{\mathrm{field}}_{m}(\mathbf w)=\mathsf{FEH}_{m}(\mathbf w),\\[2mm]
    \mathsf{QSH}^{\mathrm{bdry}}_{m}(\mathbf w), & \text{if }\mathsf{Diag}^{\mathrm{field}}_{m}(\mathbf w)=\mathsf{FEB}_{m}(\mathbf w).
  \end{cases}
\]
Define moreover the associated \emph{decoder word}
\[
  \mathsf{Word}^{\mathrm{dec}}_{m}(\mathbf w)
  \in
  \{\mathsf{HR},\mathsf{BR}\}
\]
by
\[
  \mathsf{Word}^{\mathrm{dec}}_{m}(\mathbf w)
  =
  \begin{cases}
    \mathsf{HR}, & \text{if }\mathsf{Shadow}^{\mathrm{QAM}}_{m}(\mathbf w)=\mathsf{QSH}^{\mathrm{hor}}_{m}(\mathbf w),\\[2mm]
    \mathsf{BR}, & \text{if }\mathsf{Shadow}^{\mathrm{QAM}}_{m}(\mathbf w)=\mathsf{QSH}^{\mathrm{bdry}}_{m}(\mathbf w).
  \end{cases}
\]
Interpretation: the QAM shadow is the finite state object that a later QAM-, decoder-, compression-, or router-layer is allowed to consume; the decoder word is the shortest faithful tag that still remembers whether the inherited certificate is horizon-routed or boundary-routed.
\end{definition}

\begin{proposition}[The QAM-/decoder-side shadow is a faithful finite downstream image of the certificate-aware interface]
\label{prop:qam-decoder-side-shadow-is-a-faithful-finite-downstream-image-of-the-certificate-aware-interface}
Assume the setting of Definition~\ref{def:qam-decoder-side-shadow-state}. Then:
\begin{enumerate}[label=(\roman*)]
  \item the shadow map \(\mathsf{Shadow}^{\mathrm{QAM}}_{m}\) is well-defined on every certified finite extension word;
  \item the shadow state remembers exactly the inherited field tag, forced upgrade route, and route-compatible next-work menu, but does not add new theorematic authority;
  \item the decoder word \(\mathsf{Word}^{\mathrm{dec}}_{m}(\mathbf w)\in\{\mathsf{HR},\mathsf{BR}\}\) is a faithful two-symbol readout of whether the certified state is horizon-routed or boundary-routed.
\end{enumerate}
\end{proposition}

\begin{proof}
By Proposition~\ref{prop:inherited-field-diagnostic-tag-forces-unique-conservative-upgrade-route} and Proposition~\ref{prop:every-certified-word-carries-a-finite-route-compatible-follow-on-upgrade-menu}, every certified finite extension word inherits exactly one field tag, one forced route, and one route-compatible next-work menu. Definition~\ref{def:qam-decoder-side-shadow-state} therefore assigns exactly one of the two shadow states, proving (i). Statement (ii) follows because the shadow state is merely an ordered packaging of already inherited data; it compresses the certified status but does not certify anything beyond it. For (iii), the decoder word forgets the internal menu details while preserving the decisive distinction between horizon-routed and boundary-routed certified status. Hence it is a faithful two-symbol downstream readout of the certificate-aware interface.
\end{proof}

\begin{corollary}[The present \(\sigma_3\) package already induces a first explicit QAM-/decoder-side shadow]
\label{cor:present-sigma3-package-already-induces-a-first-explicit-qam-decoder-side-shadow}
In the current active line, the first-stage word \(\mathbf w_3\in\{\mathsf A,\mathsf B\}\) already induces both a finite QAM-shadow state
\[
  \mathsf{Shadow}^{\mathrm{QAM}}_{2}(\mathbf w_3)
  \in
  \{\mathsf{QSH}^{\mathrm{hor}}_{2}(\mathbf w_3),\mathsf{QSH}^{\mathrm{bdry}}_{2}(\mathbf w_3)\}
\]
and a faithful decoder word
\[
  \mathsf{Word}^{\mathrm{dec}}_{2}(\mathbf w_3)
  \in
  \{\mathsf{HR},\mathsf{BR}\}.
\]
Thus the new extension tool layer is now explicitly attached back to the QAM-/decoder side as a finite downstream interface.
\end{corollary}

\begin{proof}
Immediate from Corollary~\ref{cor:present-sigma3-package-already-carries-a-finite-next-work-menu} and Proposition~\ref{prop:qam-decoder-side-shadow-is-a-faithful-finite-downstream-image-of-the-certificate-aware-interface}.
\end{proof}

\begin{remark}[What this QAM-/decoder-side shadow now buys]
\label{rem:what-this-qam-decoder-side-shadow-now-buys}
This is the missing explicit reattachment of the new tool layer to the QAM line. The active extension architecture can now be consumed as a finite state object instead of as a long theorematic prose chain. That is exactly what a later QAM, decoder, compression, or router implementation should see: not a new ontology, but a faithful shadow of the already certified status. In particular, the decoder word may be used for compact routing, logging, or interface normalization, but it may never be mistaken for an independent proof of continuation or boundary exposure.
\end{remark}

\noindent\textbf{Reader cue.} Read the next block as the explicit Monte-Carlo/probe discipline sitting below the certificate-aware routing and menu interface. Its purpose is to give Monte-Carlo a precise and honest role inside the current architecture: useful as a downstream structural probe and deviation profiler, but never as a substitute for theorematic certification.

\subsection*{19AZ$^{(6ai)}$. Certificate-aware Monte-Carlo / probe discipline below the routing interface}
\label{sec:certificate-aware-monte-carlo-probe-discipline-below-the-routing-interface}
After 19AZ$^{(6ag)}$ and 19AZ$^{(6ah)}$, the active line has enough structure to place Monte-Carlo and related probe machinery exactly where it belongs. The point is not to diminish that machinery, but to protect it from doing the wrong job. A good Monte-Carlo run may expose deviation patterns, stop-sensitive fingerprints, boundary clustering, transport-instability signals, or decoder-compressible distributions. What it may not do is certify a new horizon extension or erase a certified boundary stop. The present block fixes that discipline explicitly.

\begin{definition}[Certificate-aware Monte-Carlo / probe family]
\label{def:certificate-aware-monte-carlo-probe-family}
Fix a current certified horizon \(\mathbf K_m^{\mathrm{cur}}\), let \(\mathbf w\in \mathcal L_{\mathrm{ext}}^{\mathrm{term}}\) be a certified finite extension word, let
\[
  \mathsf{Shadow}^{\mathrm{QAM}}_{m}(\mathbf w)
  \in
  \{\mathsf{QSH}^{\mathrm{hor}}_{m}(\mathbf w),\mathsf{QSH}^{\mathrm{bdry}}_{m}(\mathbf w)\}
\]
be the inherited QAM-shadow state from Definition~\ref{def:qam-decoder-side-shadow-state}, and let
\[
  \mathsf{Route}^{\mathrm{upg}}_{m}(\mathbf w)
  \in
  \{\mathcal U^{\mathrm{hor}}_{m}(\mathbf w),\mathcal U^{\mathrm{bdry}}_{m}(\mathbf w)\}
\]
be the inherited upgrade route.

A \emph{certificate-aware Monte-Carlo / probe family} for \((m,\mathbf w)\) is any downstream probe procedure
\[
  \mathsf{MCProbe}_{m,\mathbf w}
\]
that takes as input only already certified data associated with \(\mathbf w\) and outputs a finite diagnostic profile
\[
  \mathsf{Prof}^{\mathrm{MC}}_{m}(\mathbf w)
\]
subject to the following conditions:
\begin{enumerate}[label=(\roman*)]
  \item \textbf{certificate inheritance:} the probe must inherit the existing shadow state and route rather than assign a new one;
  \item \textbf{route fidelity:} on the horizon route it may test only transport-side stability, continuation-side sensitivity, or compression behavior on the already certified range, while on the boundary route it may test only earliest-boundary localization, stop-sensitive clustering, defect/exposure fingerprints, or boundary-side compression behavior;
  \item \textbf{no theorematic promotion:} the output profile \(\mathsf{Prof}^{\mathrm{MC}}_{m}(\mathbf w)\) may suggest, rank, cluster, hash, Markov-compress, or Kasiski-scan deviation patterns, but it may not itself certify a later horizon \(\mathbf K_{m+1}^{\mathrm{cur}}\) or replace an inherited boundary witness \(\mathfrak B_{m+1}^{\mathrm{cur}}\);
  \item \textbf{stop discipline:} if the inherited route is boundary-routed, the probe may sharpen the earliest-boundary picture but may not sample as though theorematic continuation beyond that certified stop had already been proved.
\end{enumerate}
A probe family that violates one of these conditions is called a \emph{certificate-blind Monte-Carlo / probe family}.
\end{definition}

\begin{proposition}[Certificate-aware Monte-Carlo probes have diagnostic force but no theorematic authority]
\label{prop:certificate-aware-monte-carlo-probes-have-diagnostic-force-but-no-theorematic-authority}
Assume the setting of Definition~\ref{def:certificate-aware-monte-carlo-probe-family}. Then:
\begin{enumerate}[label=(\roman*)]
  \item every certificate-aware Monte-Carlo / probe family \(\mathsf{MCProbe}_{m,\mathbf w}\) is a legitimate downstream consumer of the inherited certificate-aware interface;
  \item the resulting profile \(\mathsf{Prof}^{\mathrm{MC}}_{m}(\mathbf w)\) may refine localization, stability, clustering, or compression questions inside the already certified route, but it does not change the theorematic status of \(\mathbf w\);
  \item any probe family that uses its output to claim a new certified extension, to smooth away an inherited boundary stop, or to inject an uncertified defect/source event is certificate-blind and therefore inadmissible in the present line.
\end{enumerate}
\end{proposition}

\begin{proof}
Clause (i) follows directly from Definition~\ref{def:certificate-aware-monte-carlo-probe-family}: the probe takes only already certified data as input and acts strictly below the inherited route. For (ii), the profile output is diagnostic by construction. It may rank or compress possibilities, but it never changes the inherited field tag, route, or next-work menu, so it cannot alter the theorematic status of \(\mathbf w\). For (iii), if a probe profile is treated as though it certified a new horizon, then it violates no-theorematic-promotion. If it smooths away a boundary-certified stop, it violates stop discipline and route fidelity. If it injects an uncertified defect/source event on a horizon-routed state, it violates route fidelity on the other side. Hence every such misuse is certificate-blind and inadmissible in the present architecture.
\end{proof}

\begin{corollary}[The present \(\sigma_3\) package already admits a first honest Monte-Carlo role]
\label{cor:present-sigma3-package-already-admits-a-first-honest-monte-carlo-role}
In the current active line, the first-stage word \(\mathbf w_3\in\{\mathsf A,\mathsf B\}\) already admits a first honest Monte-Carlo / probe role below the inherited route:
\begin{enumerate}[label=(\roman*)]
  \item if \(\mathbf w_3=\mathsf A\), Monte-Carlo may act only as a continuation-side stability or compression probe on the currently certified range;
  \item if \(\mathbf w_3=\mathsf B\), Monte-Carlo may act as an earliest-boundary localization, stop-sensitive deviation, or defect/exposure profiling probe, but not as a proof that the corridor has already continued past the certified stop.
\end{enumerate}
Thus Monte-Carlo is now explicitly fixed as a structural probe layer rather than an alternative theorematic engine.
\end{corollary}

\begin{proof}
Immediate from Corollary~\ref{cor:present-sigma3-package-already-induces-a-first-explicit-qam-decoder-side-shadow} and Proposition~\ref{prop:certificate-aware-monte-carlo-probes-have-diagnostic-force-but-no-theorematic-authority}.
\end{proof}

\begin{remark}[What this Monte-Carlo / probe discipline now buys]
\label{rem:what-this-monte-carlo-probe-discipline-now-buys}
This is the cleanest place yet to keep the mathematics-first architecture and still use aggressive probe machinery productively. Monte-Carlo may now be used to expose deviation landscapes, boundary clustering, compression-friendly fingerprints, or route-sensitive instability patterns. Those outputs can be helpful for routing attention, designing later diagnostics, or compressing large probe logs into decoder-facing summaries. But the proof burden remains exactly where it belongs: on the certified local extension line. That makes Monte-Carlo valuable without letting it impersonate theorematic closure.
\end{remark}


\noindent\textbf{Reader cue.} Read the next block as an inherited freeze audit from an earlier minor public step. Its purpose remains deliberately modest: not to claim a new global theorem, but to record in one place what the release line has stabilized, what must remain local/route-certified in the public wording, and why the present DOI-bearing freeze is honest.

\subsection*{19AZ$^{(6aj)}$. Minor-release freeze audit and public record}
\label{sec:minor-release-preparation-audit-and-freeze-checklist}
After 19AZ$^{(6x)}$--19AZ$^{(6ai)}$, the active line has turned the first post-\(\mathbf K_2^{\mathrm{cur}}\) extension step into a reusable certificate-aware tool layer. For that inherited minor-release stage, the honest question was therefore no longer whether the new line had mathematical content, but whether its public wording, routing discipline, archival preservation, and metadata had been synchronized well enough that the public freeze did not misstate what had actually been achieved. The present block records that completed audit explicitly.

\begin{definition}[Minor-release readiness audit package]
\label{def:minor-release-readiness-audit-package}
For that inherited public-release line, define the \emph{minor-release readiness audit package}
\[
  \mathsf{Audit}^{\mathrm{minor}}
  :=
  \bigl(\mathsf{Math},\mathsf{Route},\mathsf{QAM},\mathsf{MC},\mathsf{Hist},\mathsf{Meta}\bigr)
\]
as follows:
\begin{enumerate}[label=(\roman*)]
  \item \(\mathsf{Math}\): the live OP~1 kernel is discharged locally on the retained no-second-carrier / strict-upper-gap corridor, and the live OP~5 Step~(B) kernel is discharged locally on the retained branch-shell / shell-transfer corridor;
  \item \(\mathsf{Route}\): the first later leading admissible slot beyond the current certified horizon is normalized by a finite decision datum, an action interface, an admissible terminal word language, field tags, route tags, and a finite next-work menu;
  \item \(\mathsf{QAM}\): the new extension layer is attached back to the QAM/decoder side by a finite shadow-state interface and the faithful decoder words \(\{\mathsf{HR},\mathsf{BR}\}\);
  \item \(\mathsf{MC}\): Monte-Carlo and related probe machinery are explicitly confined to a certificate-aware downstream diagnostic role;
  \item \(\mathsf{Hist}\): the full historical witness remains preserved in the monolithic source as Part~II rather than as a shortened fragment;
  \item \(\mathsf{Meta}\): title-line DOI logic, version/chronology wording, and release wording are synchronized so that the concept DOI remains the stable title-page pointer while public version-specific DOI claims appear only in the chronology or explicit release notes.
\end{enumerate}
\end{definition}

\begin{proposition}[Present release line is mathematically minor-ready and release-packaged]
\label{prop:present-release-line-is-mathematically-minor-ready-and-release-packaged}
Assume the audit package from Definition~\ref{def:minor-release-readiness-audit-package}. Then:
\begin{enumerate}[label=(\roman*)]
  \item the mathematics-side burden of that inherited minor release is strong enough for an honest public freeze: the new contribution is a certificate-aware extension/routing/probe layer above the public v3.28.0 infrastructure together with local retained-corridor closure of the live OP~1 and OP~5 Step~(B) kernels;
  \item the public wording of that freeze remains local and route-certified, i.e. it reports retained-corridor kernel closure and the new certificate-aware tool layer, but it does not flatten those results into a new unconditional global closure slogan;
  \item the release-packaging tasks for that inherited public minor release have been synchronized: version/date metadata, DOI logic, historical preservation, and public wording have been normalized, while ordinary citation/hyperref cleanup can continue in later working states without changing the mathematics.
\end{enumerate}
\end{proposition}

\begin{proof}
Clause (i) is exactly the content accumulated in 19AZ$^{(6x)}$--19AZ$^{(6ai)}$: finite next-slot decision data, faithful output tokens, action/word routing, certificate-aware matrix/crystal/field readout, admissibility and routing control, finite next-work menus, finite QAM shadows, and an honest downstream Monte-Carlo role. These are not merely editorial restatements but reusable structural tools. The local retained-corridor OP~1 and OP~5 Step~(B) closure statements are already recorded in the active main line, so the mathematics side no longer depends on introducing yet another theorematic layer before that inherited minor public freeze. Clause (ii) follows from the route-certified architecture itself: every new statement in this block is built from inherited certified corridors, finite decision packaging, and certificate-aware downstream consumers. Flattening that into an unconditional global slogan would destroy the very distinction the new layer was designed to preserve. Clause (iii) follows because what is left to do does not add new mathematical content; it concerns release hygiene, metadata synchronization, and wording discipline.
\end{proof}

\begin{corollary}[Honest public release message for that inherited minor release]
\label{cor:honest-public-release-message-for-the-present-minor-release}
For that inherited public minor release, the honest release message is the following:
\begin{quote}\small
The public v3.28.0 line is carried forward into the public v\docversion\ release by a certificate-aware extension/routing/probe layer. On top of that inherited infrastructure, the live OP~1 kernel and the live OP~5 Step~(B) kernel are closed locally on the standing retained corridors. The resulting public gain is therefore not a new global closure slogan, but a route-certified tool layer together with retained-corridor kernel closure, finite decoder/QAM shadows, and a disciplined downstream role for Monte-Carlo and related probes.
\end{quote}
\end{corollary}

\begin{proof}
Immediate from Proposition~\ref{prop:present-release-line-is-mathematically-minor-ready-and-release-packaged}.
\end{proof}

\begin{remark}[Freeze checklist completed for the DOI-bearing release]
\label{rem:freeze-checklist-completed-for-the-doi-bearing-release}
In that inherited freeze block, the deliberately boring tasks had been carried out: version/date metadata in the wrapper had been synchronized, chronology wording recorded v3.28.0 as the prior public baseline and the then-current line as the public v\docversion\ release, the historical witness remained preserved in full, and the front-matter wording had been normalized so that local retained-corridor results were not overstated as unconditional global closure. Ordinary citation/hyperref cleanup could still improve future working states, but it no longer blocked the DOI-bearing public freeze.
\end{remark}


\noindent\textbf{Reader cue.} Read the next block as the first post-release working bridge that couples the certificate-aware upgrade menus to honest probe consumption on the matrix, crystal, and field-diagnostic side. Its purpose is not to reopen the theorematic routing decision, but to package conservative sharpening together with route-compatible probe use into one downstream working object.

